% Complete reviewed Unit 18; editorial master: chapters/ch18/answer-key.md
% Markdown SHA-256: 6370dfd1a9d91d2bf0539d469d710976a88125fa5716d716e552bd4e93aeb033
% One-use initial assembly followed by complete content and page review.
\newcommand{\runningtitle}{Unit 18. Answer key}
\input{../book-layout.tex}
\hyphenpenalty=10000
\exhyphenpenalty=10000
\emergencystretch=3em
\newcommand{\itempara}[2]{\par\addvspace{3pt}\noindent\hangindent=1.8em\hangafter=1\makebox[1.8em][l]{#1.}#2\par}
\newcommand{\reserve}[1]{\par\penalty-100\begingroup\dimen0=\pagegoal\advance\dimen0 by-\pagetotal\ifdim\dimen0<#1\newpage\fi\endgroup}
\begin{document}
\addfontfeatures{Ligatures=TeXOff}
\raggedright
\booktitle{UNIT 18}{Answer key}

% mdparagraph:P001:begin
The composition models are new teaching material. They retain Cook’s complete Exercises 44–48 under his joint group \textbf{XII, XIII}. The connected reading is the separate Apollonius passage, and its translation was prepared independently from the actual Old English before comparison with other complete translations. Each model is one defensible answer; a different grammatical answer can still miss an explicit lexical, mood, number, omission or order target.\textsuperscript{1}
% mdparagraph:P001:end

% mdparagraph:P002:begin
Preserve every participant, action, time relation and clause connection. Except where an order is specified, other grammatical orders are acceptable. English translations and analyses may use equivalent precise wording. Documented variants, such as \textbf{ðām} beside selected \textbf{ðǣm}, are not automatically errors; tasks that expressly require the selected forms still test those forms. The notes below identify the distinctions most likely to change the meaning.
% mdparagraph:P002:end

% mdparagraph:P003:begin
The original hints and complete local help remain available. A copied supplied form is not independent recall of an unassigned paradigm. In particular, source names, quantities, rare reading frames and the combined counterfactual in 44.7 are used with their help visible. A simple preterite may render English “did,” “has done” or “had done” where the context establishes the relation; a subjunctive in a report does not by itself make the report false.
% mdparagraph:P003:end

\reserve{7\baselineskip}
% mdparagraph:P004:begin
\section*{A. Clauses and reference}
% mdparagraph:P004:end

\reserve{8\baselineskip}
% mdparagraph:P005:begin
\subsection*{A01}
% mdparagraph:P005:end

% key: KEY-NEW-U18-A01
% mdparagraph:P006:begin
\textbf{Task:} Analyse the two sentences. No further context identifies the reader.
% mdparagraph:P006:end

% mdparagraph:P007:begin
\textbf{Input:} Se cyning sealde ðǣm męnn ðā bōc. Hē hīe rǣdde.
% mdparagraph:P007:end

% output: NEW-U18-A01-O001
\reserve{5\baselineskip}
% mdparagraph:P008:begin
\textbf{Identify the giver, recipient and thing given.}
% mdparagraph:P008:end

% mdparagraph:P009:begin
The king gave the book to the man. \textbf{Se cyning} is the nominative giver, \textbf{ðǣm męnn} the dative singular recipient, and \textbf{ðā bōc} the accusative singular gift. The article and giving frame identify męnn as singular dative here; its resemblance to a plural form does not make several recipients.
% mdparagraph:P009:end

% output: NEW-U18-A01-O002
\reserve{5\baselineskip}
% mdparagraph:P010:begin
\textbf{Give the referent and case of hīe.}
% mdparagraph:P010:end

% mdparagraph:P011:begin
\textbf{Hīe} is feminine accusative singular, referring to \textbf{bōc}, the book read. Grammatical gender allows the same pronoun to refer to an inanimate feminine noun; “it” is the natural English rendering. Reading it as “they” would also conflict with its object role here.
% mdparagraph:P011:end

% output: NEW-U18-A01-O003
\reserve{5\baselineskip}
% mdparagraph:P012:begin
\textbf{Which named people can hē refer to? Explain what the forms establish.}
% mdparagraph:P012:end

% mdparagraph:P013:begin
Either the king or the man can be the reader. Both are masculine singular and can be resumed by nominative \textbf{hē}, the subject of \textbf{rǣdde}. The stated absence of further context leaves that choice unresolved. Neither proximity to męnn nor continuity with the earlier subject is a grammatical requirement.
% mdparagraph:P013:end

% output: NEW-U18-A01-O004
\reserve{5\baselineskip}
\par\begin{minipage}{\linewidth}\raggedright
% mdparagraph:P014:begin
\textbf{Rewrite only the second sentence so that the recipient is explicitly the reader.}
% mdparagraph:P014:end

% mdparagraph:P015:begin
\textbf{Se mann hīe rǣdde.}
% mdparagraph:P015:end
\end{minipage}\par

% mdparagraph:P016:begin
The man is now the explicitly named nominative subject; \textbf{hīe} still refers to the book as the accusative object. The result means “The man read it.” A grammatical rearrangement may preserve those roles, but retaining hē alone leaves the ambiguity that the task asks you to remove.
% mdparagraph:P016:end

\reserve{8\baselineskip}
% mdparagraph:P017:begin
\subsection*{A02}
% mdparagraph:P017:end

% key: KEY-NEW-U18-A02
% mdparagraph:P018:begin
\textbf{Task:} Convert this direct statement into a report. Retain the king as the person who saw the man; use sęcgan with a ðæt clause and a preterite subjunctive in that clause.
% mdparagraph:P018:end

% mdparagraph:P019:begin
\textbf{Input:} Se cyning cwæð: “Ic geseah ðone mann.”
% mdparagraph:P019:end

% mdparagraph:P020:begin
\textbf{Input:} sęcgan
% mdparagraph:P020:end

% mdparagraph:P021:begin
\textbf{Input:} ðæt
% mdparagraph:P021:end

% output: NEW-U18-A02-O001
\reserve{5\baselineskip}
% mdparagraph:P022:begin
\textbf{Identify the referent of ic and the time of the seeing relative to the quoted statement.}
% mdparagraph:P022:end

% mdparagraph:P023:begin
\textbf{Ic} is the king, the speaker inside the quotation. \textbf{Geseah} presents his seeing as earlier than his quoted statement: “I saw the man.” It does not refer to the narrator who introduces the quotation.
% mdparagraph:P023:end

% output: NEW-U18-A02-O002
\reserve{5\baselineskip}
\par\begin{minipage}{\linewidth}\raggedright
% mdparagraph:P024:begin
\textbf{Write the complete reported statement.}
% mdparagraph:P024:end

% mdparagraph:P025:begin
\textbf{Se cyning sægde ðæt hē ðone mann gesāwe.}
% mdparagraph:P025:end
\end{minipage}\par

% mdparagraph:P026:begin
“The king said that he had seen the man.” \textbf{Sægde} introduces reported content; \textbf{hē} keeps the king as the seer, and \textbf{gesāwe} is preterite subjunctive third singular. \textbf{Ðone mann} remains the accusative person seen. Repeating se cyning inside the clause can make the reference explicit. An indicative report can be valid language, but does not meet this task’s requested subjunctive.
% mdparagraph:P026:end

% output: NEW-U18-A02-O003
\reserve{5\baselineskip}
% mdparagraph:P027:begin
\textbf{Explain whether the chosen subjunctive proves that the king’s statement is false.}
% mdparagraph:P027:end

% mdparagraph:P028:begin
No. Here the subjunctive presents the reported content and does not establish its truth or falsity. The seeing is earlier because of the original statement’s temporal viewpoint, not because every preterite subjunctive must express an earlier event or a false proposition.
% mdparagraph:P028:end

\reserve{8\baselineskip}
% mdparagraph:P029:begin
\subsection*{A03}
% mdparagraph:P029:end

% key: KEY-NEW-U18-A03
% mdparagraph:P030:begin
\textbf{Task:} Analyse the sentence, then replace the men by the man and keep the king singular.
% mdparagraph:P030:end

% mdparagraph:P031:begin
\textbf{Input:} Ðā męnn ðe ðone cyning gesāwon, sungon.
% mdparagraph:P031:end

% output: NEW-U18-A03-O001
\reserve{5\baselineskip}
% mdparagraph:P032:begin
\textbf{Translate the complete sentence.}
% mdparagraph:P032:end

% mdparagraph:P033:begin
“The men who saw the king sang.” “The men that saw the king sang” has the same meaning. The men perform both actions; the king is the person seen, not a seer or the singer.
% mdparagraph:P033:end

% output: NEW-U18-A03-O002
\reserve{5\baselineskip}
% mdparagraph:P034:begin
\textbf{Identify the role represented by ðe and explain the agreement of both finite verbs.}
% mdparagraph:P034:end

% mdparagraph:P035:begin
\textbf{Ðe} represents the men as the subject of \textbf{gesāwon}, preterite indicative third plural. The main subject \textbf{ðā męnn} governs \textbf{sungon}, also preterite indicative third plural. \textbf{Ðone cyning} is accusative singular inside the relative clause and controls neither verb’s number.
% mdparagraph:P035:end

% output: NEW-U18-A03-O003
\reserve{5\baselineskip}
\par\begin{minipage}{\linewidth}\raggedright
% mdparagraph:P036:begin
\textbf{Write the complete singular version.}
% mdparagraph:P036:end

% mdparagraph:P037:begin
\textbf{Se mann ðe ðone cyning geseah, sang.}
% mdparagraph:P037:end
\end{minipage}\par

% mdparagraph:P038:begin
“The man who saw the king sang.” \textbf{Se mann} replaces the plural subject; \textbf{geseah} and \textbf{sang} are third singular. The object \textbf{ðone cyning} stays singular. Changing only one finite verb would leave the other action attached to an unrequested plural subject.
% mdparagraph:P038:end

\reserve{7\baselineskip}
% mdparagraph:P039:begin
\section*{C. Composition}
% mdparagraph:P039:end

\reserve{8\baselineskip}
% mdparagraph:P040:begin
\subsection*{C01}
% mdparagraph:P040:end

% key: KEY-NEW-U18-C01
% mdparagraph:P041:begin
\textbf{Task:} Translate: The man whom the king saw gave the servant the book. Use ðe and sęllan.
% mdparagraph:P041:end

% mdparagraph:P042:begin
\textbf{Input:} ðe
% mdparagraph:P042:end

% mdparagraph:P043:begin
\textbf{Input:} sęllan
% mdparagraph:P043:end

% output: NEW-U18-C01-O001
\reserve{5\baselineskip}
\par\begin{minipage}{\linewidth}\raggedright
% mdparagraph:P044:begin
\textbf{Write one complete Old English sentence.}
% mdparagraph:P044:end

% mdparagraph:P045:begin
\textbf{Se mann ðe se cyning geseah, sealde ðǣm cnihte ðā bōc.}
% mdparagraph:P045:end
\end{minipage}\par

% mdparagraph:P046:begin
The man is the main nominative subject of \textbf{sealde}. Inside the relative, \textbf{se cyning} is the subject of \textbf{geseah}, and invariant \textbf{ðe} represents the man as its accusative object. \textbf{Ðǣm cnihte} is the servant receiving the accusative book. Word order may vary while those roles remain clear; making the king give the book changes the meaning.
% mdparagraph:P046:end

\reserve{8\baselineskip}
% mdparagraph:P047:begin
\subsection*{C02}
% mdparagraph:P047:end

% key: KEY-NEW-U18-C02
% mdparagraph:P048:begin
\textbf{Task:} Translate: The king ordered the men to give the servant the books. Use hātan and a plain infinitive of sęllan.
% mdparagraph:P048:end

% mdparagraph:P049:begin
\textbf{Input:} hātan
% mdparagraph:P049:end

% mdparagraph:P050:begin
\textbf{Input:} sęllan
% mdparagraph:P050:end

% output: NEW-U18-C02-O001
\reserve{5\baselineskip}
\par\begin{minipage}{\linewidth}\raggedright
% mdparagraph:P051:begin
\textbf{Write one complete Old English sentence.}
% mdparagraph:P051:end

% mdparagraph:P052:begin
\textbf{Se cyning hēt ðā męnn ðǣm cnihte ðā bēc sęllan.}
% mdparagraph:P052:end
\end{minipage}\par

% mdparagraph:P053:begin
The king commanded the men to give the servant the books. \textbf{Hēt} has accusative \textbf{ðā męnn} as the instructed actors of plain infinitive \textbf{sęllan}. The servant is the dative recipient and \textbf{ðā bēc} the accusative plural gift. Although męnn has the same spelling in nominative and accusative plural, it is accusative in this construction. A finite plural verb in place of sęllan misses the requested infinitive frame.
% mdparagraph:P053:end

\clearpage
% mdparagraph:P054:begin
\section*{F. Forms and grammatical roles}
% mdparagraph:P054:end

\reserve{8\baselineskip}
% mdparagraph:P055:begin
\subsection*{F01}
% mdparagraph:P055:end

% key: KEY-NEW-U18-F01
% mdparagraph:P056:begin
\textbf{Task:} Translate each sentence, then explain the selected ending contrast. Use the newly taught words for ass, sack and wheat.
% mdparagraph:P056:end

% output: NEW-U18-F01-O001
\reserve{5\baselineskip}
\par\begin{minipage}{\linewidth}\raggedright
% mdparagraph:P057:begin
\textbf{The ass has the sack.}
% mdparagraph:P057:end

% mdparagraph:P058:begin
\textbf{Se assa hæfð ðone sacc.}
% mdparagraph:P058:end
\end{minipage}\par

% mdparagraph:P059:begin
\textbf{Se assa} is nominative singular; \textbf{ðone sacc} is the accusative singular possession. \textbf{Hæfð} agrees with the one ass. The weak noun’s nominative assa must not be replaced by oblique singular assan.
% mdparagraph:P059:end

% output: NEW-U18-F01-O002
\reserve{5\baselineskip}
\par\begin{minipage}{\linewidth}\raggedright
% mdparagraph:P060:begin
\textbf{The asses have the sacks.}
% mdparagraph:P060:end

% mdparagraph:P061:begin
\textbf{Ðā assan habbað ðā saccas.}
% mdparagraph:P061:end
\end{minipage}\par

% mdparagraph:P062:begin
\textbf{Ðā assan} is weak nominative plural; \textbf{ðā saccas} is strong accusative plural. \textbf{Habbað} agrees with the plural possessors. The requested two noun patterns differ: an invented strong plural of assa would miss its paradigm.
% mdparagraph:P062:end

% output: NEW-U18-F01-O003
\reserve{5\baselineskip}
\par\begin{minipage}{\linewidth}\raggedright
% mdparagraph:P063:begin
\textbf{The men give wheat to the asses.}
% mdparagraph:P063:end

% mdparagraph:P064:begin
\textbf{Ðā męnn sęllað ðǣm assum hwǣte.}
% mdparagraph:P064:end
\end{minipage}\par

% mdparagraph:P065:begin
The men are the nominative plural givers and govern \textbf{sęllað}. \textbf{Ðǣm assum} is dative plural, the recipients; \textbf{hwǣte} is the accusative mass singular thing given. Grammatical changes of object order are possible. Making the asses the things given, or inventing a count plural of wheat, changes the required relations.
% mdparagraph:P065:end

% output: NEW-U18-F01-O004
\reserve{5\baselineskip}
\par\begin{minipage}{\linewidth}\raggedright
% mdparagraph:P066:begin
\textbf{The wheat is in the sacks.}
% mdparagraph:P066:end

% mdparagraph:P067:begin
\textbf{Se hwǣte is on ðǣm saccum.}
% mdparagraph:P067:end
\end{minipage}\par

% mdparagraph:P068:begin
The mass singular \textbf{se hwǣte} governs \textbf{is}. \textbf{On ðǣm saccum} gives a static location, hence dative plural. The several sacks do not turn the wheat into a plural subject or require a plural copula.
% mdparagraph:P068:end

% output: NEW-U18-F01-O005
\reserve{5\baselineskip}
% mdparagraph:P069:begin
\textbf{Compare the plural noun endings in the first three translations and the location phrase.}
% mdparagraph:P069:end

% mdparagraph:P070:begin
In the plural sentence, \textbf{assan} has the weak nominative plural ending \textbf{-an}, whereas \textbf{saccas} has strong accusative plural \textbf{-as}. Both classes have dative plural \textbf{-um}: \textbf{assum} as recipients and \textbf{saccum} after static on. Wheat remains mass singular. A shared dative ending does not make all the forms of the two declensions identical.
% mdparagraph:P070:end

\reserve{8\baselineskip}
% mdparagraph:P071:begin
\subsection*{F02}
% mdparagraph:P071:end

% key: KEY-NEW-U18-F02
% mdparagraph:P072:begin
\textbf{Task:} Translate the two sentences using the newly taught verb meaning need. Then explain the governed noun.
% mdparagraph:P072:end

% output: NEW-U18-F02-O001
\reserve{5\baselineskip}
\par\begin{minipage}{\linewidth}\raggedright
% mdparagraph:P073:begin
\textbf{The man needs wheat.}
% mdparagraph:P073:end

% mdparagraph:P074:begin
\textbf{Se mann beðearf hwǣtes.}
% mdparagraph:P074:end
\end{minipage}\par

% mdparagraph:P075:begin
\textbf{Beðearf} is present indicative third singular, agreeing with the man. The selected need frame governs genitive singular \textbf{hwǣtes}. English places “wheat” directly after “needs,” but that does not determine its Old English case.
% mdparagraph:P075:end

% output: NEW-U18-F02-O002
\reserve{5\baselineskip}
\par\begin{minipage}{\linewidth}\raggedright
% mdparagraph:P076:begin
\textbf{The men needed wheat.}
% mdparagraph:P076:end

% mdparagraph:P077:begin
\textbf{Ðā męnn beðorfton hwǣtes.}
% mdparagraph:P077:end
\end{minipage}\par

% mdparagraph:P078:begin
\textbf{Beðorfton} is preterite indicative third plural: the men needed wheat. \textbf{Hwǣtes} remains genitive mass singular. Present plural \textbf{beðurfon} would give the right subject number but the wrong time.
% mdparagraph:P078:end

% output: NEW-U18-F02-O003
\reserve{5\baselineskip}
% mdparagraph:P079:begin
\textbf{Identify the tense and number of each finite verb, and explain the case of wheat in both sentences.}
% mdparagraph:P079:end

% mdparagraph:P080:begin
\textbf{Beðearf} is present indicative third singular; \textbf{beðorfton} is preterite indicative third plural. \textbf{Hwǣtes} is genitive singular in both because beðurfan governs the thing needed in this selected frame. Changing the subject’s number does not change the mass noun’s number or its government. Other attested government does not meet this exercise’s genitive target.
% mdparagraph:P080:end

\reserve{8\baselineskip}
% mdparagraph:P081:begin
\subsection*{F03}
% mdparagraph:P081:end

% key: KEY-NEW-U18-F03
% mdparagraph:P082:begin
\textbf{Task:} Translate the four sentences using the newly taught verbs meaning recognize and undo/\allowbreak{}open. Then explain the roles and number changes.
% mdparagraph:P082:end

% output: NEW-U18-F03-O001
\reserve{5\baselineskip}
\par\begin{minipage}{\linewidth}\raggedright
% mdparagraph:P083:begin
\textbf{I recognize the king.}
% mdparagraph:P083:end

% mdparagraph:P084:begin
\textbf{Ic oncnāwe ðone cyning.}
% mdparagraph:P084:end
\end{minipage}\par

% mdparagraph:P085:begin
\textbf{Oncnāwe} is present indicative first singular. \textbf{Ðone cyning} is the singular accusative person recognized. The long ā belongs to this present stem; a past form would change “recognize” to “recognized.”
% mdparagraph:P085:end

% output: NEW-U18-F03-O002
\reserve{5\baselineskip}
\par\begin{minipage}{\linewidth}\raggedright
% mdparagraph:P086:begin
\textbf{We recognized the king.}
% mdparagraph:P086:end

% mdparagraph:P087:begin
\textbf{Wē oncnēowon ðone cyning.}
% mdparagraph:P087:end
\end{minipage}\par

% mdparagraph:P088:begin
\textbf{Oncnēowon} is preterite indicative first plural, with the past \textbf{cnēow-} stem. Only the recognizers change number; \textbf{ðone cyning} remains one accusative person. Pluralizing the king would add a change the task does not request.
% mdparagraph:P088:end

% output: NEW-U18-F03-O003
\reserve{5\baselineskip}
\par\begin{minipage}{\linewidth}\raggedright
% mdparagraph:P089:begin
\textbf{The man opens the sack.}
% mdparagraph:P089:end

% mdparagraph:P090:begin
\textbf{Se mann undēð ðone sacc.}
% mdparagraph:P090:end
\end{minipage}\par

% mdparagraph:P091:begin
\textbf{Undēð} is present indicative third singular; the man opens the singular accusative sack. Its anomalous \textbf{dēð} stem must not be replaced by the infinitive undōn in a finite assertion.
% mdparagraph:P091:end

% output: NEW-U18-F03-O004
\reserve{5\baselineskip}
\par\begin{minipage}{\linewidth}\raggedright
% mdparagraph:P092:begin
\textbf{The men opened the sacks.}
% mdparagraph:P092:end

% mdparagraph:P093:begin
\textbf{Ðā męnn undydon ðā saccas.}
% mdparagraph:P093:end
\end{minipage}\par

% mdparagraph:P094:begin
\textbf{Undydon} is preterite indicative third plural. Both subject and object are plural: \textbf{ðā męnn} opened \textbf{ðā saccas}. The preterite \textbf{dyde-} system is distinct from present undēð; keeping the singular object or present tense would miss part of the change.
% mdparagraph:P094:end

% output: NEW-U18-F03-O005
\reserve{5\baselineskip}
% mdparagraph:P095:begin
\textbf{Identify the direct objects and explain which words change when the subject, object or tense changes.}
% mdparagraph:P095:end

% mdparagraph:P096:begin
The king is the accusative singular object in both recognition sentences. The sack and sacks are accusative objects in the opening sentences. Subject person and number determine finite agreement, and the requested time determines the verb stem; object number determines its own article and noun. The nearer object does not govern the verb’s person or number.
% mdparagraph:P096:end

\clearpage
% mdparagraph:P097:begin
\section*{Cook 44}
% mdparagraph:P097:end

\par\begin{minipage}{\linewidth}\raggedright
% mdparagraph:P098:begin
\subsection*{Cook 44.1}
% mdparagraph:P098:end

% key: KEY-EX-XII-XIII-44-01
% mdparagraph:P099:begin
Iacob cwæð tō his sunum: Gāð and bycgað ūs sumne dǣl mętes.
% mdparagraph:P099:end

% output: EX-XII-XIII-44-01-O001
% mdparagraph:P100:begin
\textbf{Cook’s prompt:} Jacob (Iacob) said to his sons, Go and buy us a certain amount of food (męte).
% mdparagraph:P100:end

% mdparagraph:P101:begin
\textbf{Target and explanation:} \textbf{Gāð} and \textbf{bycgað} are plural commands to the sons. \textbf{Ūs} is dative beneficiary, “for us”; \textbf{sumne dǣl mętes} is the accusative amount bought, with genitive mętes specifying the food. The amount phrase, not a bare nominative food noun, preserves “a certain amount.” Jacob’s quoted commands must not become a narrative that the sons already went and bought.
% mdparagraph:P101:end

\end{minipage}\par\addvspace{4pt}

\par\begin{minipage}{\linewidth}\raggedright
% mdparagraph:P102:begin
\subsection*{Cook 44.2}
% mdparagraph:P102:end

% key: KEY-EX-XII-XIII-44-02
% mdparagraph:P103:begin
Ðā cwæð Iudas: Ðæs landes ealdor swōr āðas ætforan ūs ðæt wē ne mōsten cuman ætforan him, būtan wē ūrne gingstan brōðor mid ūs bringen.
% mdparagraph:P103:end

% output: EX-XII-XIII-44-02-O001
% mdparagraph:P104:begin
\textbf{Cook’s prompt:} Then said Judah (Iudas), The chief of the land swore oaths before us that we should not (mōtan) come before him, except we bring our youngest brother with us.
% mdparagraph:P104:end

% mdparagraph:P105:begin
\textbf{Target and explanation:} Judah reports the chief’s oath. \textbf{Swōr} is past indicative; \textbf{mōsten} is first-plural past subjunctive in the reported prohibition, with plain \textbf{cuman}. It concerns permission to appear before the chief, not inability to travel. Supplied conditional \textbf{būtan} introduces the exception, and \textbf{bringen} keeps the brothers as first-plural subject. \textbf{Ūrne gingstan brōðor} is the accusative companion brought; the chief remains singular \textbf{him}, dative after \textbf{ætforan} even with \textbf{cuman}.
% mdparagraph:P105:end

% mdparagraph:P106:begin
\textbf{Supplied construction:} The learner is given \textbf{ætforan ūs} and \textbf{cuman ætforan him}, both with dative complements. The preposition determines the case; movement alone does not license \textbf{hine} here.
% mdparagraph:P106:end

\end{minipage}\par\addvspace{4pt}

\par\begin{minipage}{\linewidth}\raggedright
% mdparagraph:P107:begin
\subsection*{Cook 44.3}
% mdparagraph:P107:end

% key: KEY-EX-XII-XIII-44-03
% mdparagraph:P108:begin
Gif ðū hine mid ūs sęndan wilt, wē gāð ætgædere and bycgað ðā ðing ðe wē beðurfon.
% mdparagraph:P108:end

% output: EX-XII-XIII-44-03-O001
% mdparagraph:P109:begin
\textbf{Cook’s prompt:} If thou wilt send him with us (verb last), we (om. will) go together and buy the things that (87. d) we need (beðurfan, 131).
% mdparagraph:P109:end

% mdparagraph:P110:begin
\textbf{Target and explanation:} The stipulated final verb in the condition is \textbf{wilt}, governing plain \textbf{sęndan}. The two main presents \textbf{gāð} and \textbf{bycgað} have subject \textbf{wē} and express the proposed future actions without an added “will.” \textbf{Ðā ðing} is accusative object of buying; invariant \textbf{ðe} has the complement role required by \textbf{beðurfon} inside the relative. The selected need frame is genitive, although ðe displays no case ending. Treating both main verbs as commands changes the speakers’ undertaking.
% mdparagraph:P110:end

\end{minipage}\par\addvspace{4pt}

\par\begin{minipage}{\linewidth}\raggedright
% mdparagraph:P111:begin
\subsection*{Cook 44.4}
% mdparagraph:P111:end

% key: KEY-EX-XII-XIII-44-04
% mdparagraph:P112:begin
Gē earmlice wið mē dydon, ðā gē him sægdon ðæt gē mā brōðra hæfden.
% mdparagraph:P112:end

% output: EX-XII-XIII-44-04-O001
% mdparagraph:P113:begin
\textbf{Cook’s prompt:} Ye dealt ill (earmlice) with (wið) me, when ye said to him that ye had more (77) brethren (154).
% mdparagraph:P113:end

% mdparagraph:P114:begin
\textbf{Target and explanation:} The complete supplied phrase \textbf{earmlice wið mē dydon} means “you dealt badly with me.” Both \textbf{dydon} and \textbf{sægdon} agree with plural gē. \textbf{Hæfden} is second-plural past subjunctive in what they said; \textbf{mā brōðra} has a quantity word followed by genitive plural. It does not mean “more to the brothers,” nor does the report’s subjunctive establish that their statement was false.
% mdparagraph:P114:end

\end{minipage}\par\addvspace{4pt}

\par\begin{minipage}{\linewidth}\raggedright
% mdparagraph:P115:begin
\subsection*{Cook 44.5}
% mdparagraph:P115:end

% key: KEY-EX-XII-XIII-44-05
% mdparagraph:P116:begin
Se mann āscode ūs ealle hwæðer ūre fæder lifde and hwæðer wē brōðor hæfden, and wē him andswarodon; wē ne wiston ðæt hē wolde giernan ðæs ðæt wē ūrne brōðor ðider lǣdden.
% mdparagraph:P116:end

% output: EX-XII-XIII-44-05-O001
% mdparagraph:P117:begin
\textbf{Cook’s prompt:} The man asked (āscian) us all whether our father was living [lived] and whether we had brethren, and we answered him; we knew not that he would (willan) desire (giernan) that (156. a), that we should bring (lǣdan, 194. b) our brother thither.
% mdparagraph:P117:end

% mdparagraph:P118:begin
\textbf{Target and explanation:} \textbf{Ūs ealle} retains “us all.” The two \textbf{hwæðer} clauses are indirect questions: whether the father \textbf{lifde}, and whether the speakers \textbf{hæfden} brothers. \textbf{Wiston} is plural past knowing; \textbf{wolde} presents the chief’s contemplated desire. Retain \textbf{giernan ðæs ðæt}: genitive ðæs is the thing desired, and the following content clause explains it. \textbf{Lǣdden} is first-plural past subjunctive, with their brother as accusative object. Omitting either question, all, or the genitive demonstrative loses a source control.
% mdparagraph:P118:end

\end{minipage}\par\addvspace{4pt}

\par\begin{minipage}{\linewidth}\raggedright
% mdparagraph:P119:begin
\subsection*{Cook 44.6}
% mdparagraph:P119:end

% key: KEY-EX-XII-XIII-44-06
% mdparagraph:P120:begin
Iudas cwæð: Sęnd ðone cnapan mid mē, ðæt wē mægen bycgan ðā ðing ðe wē beðurfon, ðæt wē ne forweorðen; būtan ic hine ęft hām bringe, sīe ic scyldig.
% mdparagraph:P120:end

% output: EX-XII-XIII-44-06-O001
% mdparagraph:P121:begin
\textbf{Cook’s prompt:} Judah said, Send the boy with me, that we may (mugan) buy the things that we need, that we perish (forweorðan) not; except I bring him back home, may I be (193. a) guilty (scyldig).
% mdparagraph:P121:end

% mdparagraph:P122:begin
\textbf{Target and explanation:} Singular \textbf{sęnd} addresses Jacob. \textbf{Ðæt wē mægen bycgan} gives the purpose of sending the boy, and \textbf{ðæt wē ne forweorðen} gives the further purpose, avoiding starvation or death. The needed things retain the relative dependency explained in 44.3. In the new condition the subject changes to \textbf{ic}: \textbf{bringe} is first-singular present subjunctive, and \textbf{sīe ic scyldig} is the resulting self-imprecation, “let me be guilty.” A simple indicative admission of guilt would lose that force.
% mdparagraph:P122:end

\end{minipage}\par\addvspace{4pt}

\par\begin{minipage}{\linewidth}\raggedright
% mdparagraph:P123:begin
\subsection*{Cook 44.7}
% mdparagraph:P123:end

% key: KEY-EX-XII-XIII-44-07
% mdparagraph:P124:begin
Gif ðēos ęlcung nǣre, wē wǣren ōðre sīðe cumen.
% mdparagraph:P124:end

% output: EX-XII-XIII-44-07-O001
% mdparagraph:P125:begin
\textbf{Cook’s prompt:} Were it not for this delay [if this delay were not], we had [were, 188] come (om. a) second time (sīð, 176).
% mdparagraph:P125:end

% mdparagraph:P126:begin
\textbf{Target and explanation:} Use the entire supplied counterfactual frame. \textbf{Nǣre} is negative past subjunctive third singular with “this delay”; \textbf{wǣren} is first-plural past subjunctive auxiliary with past participle \textbf{cumen}. Thus the result is “we would have come a second time,” not merely “we were coming.” \textbf{Ōðre sīðe} is adverbial instrumental singular under Cook 176; the noun hint sīð does not remove its ending. No article represents English “a.” An ordinary accusative phrase or a finite reading of cumen misses the stipulated construction.
% mdparagraph:P126:end

\end{minipage}\par\addvspace{4pt}

\par\begin{minipage}{\linewidth}\raggedright
% mdparagraph:P127:begin
\subsection*{Cook 44.8}
% mdparagraph:P127:end

% key: KEY-EX-XII-XIII-44-08
% mdparagraph:P128:begin
Ðā cwæð Iacob: Gif gē nīede swā dōn sculon, dōð swā gē willað; nimað of ēowrum bętstum wæstmum on fatum, and bringað ðǣm męnn lāc, hunig and æcernu.
% mdparagraph:P128:end

% output: EX-XII-XIII-44-08-O001
% mdparagraph:P129:begin
\textbf{Cook’s prompt:} Then said Jacob, If ye needs (nīede, 176) must do so, do as ye will; take of (of) your best fruits in vessels, and carry (bringan) presents (om. to) the man, honey (hunig, 47) and acorns (æcern, 47).
% mdparagraph:P129:end

% mdparagraph:P130:begin
\textbf{Target and explanation:} \textbf{Nīede} supplies necessity, while \textbf{sculon} governs dōn. The following \textbf{dōð, nimað, bringað} are plural commands. \textbf{Of ēowrum bętstum wæstmum} and \textbf{on fatum} use dative plurals; \textbf{ðǣm męnn} is the dative recipient without an added tō. \textbf{Lāc} is accusative plural “presents,” despite matching its singular spelling; honey is mass singular and \textbf{æcernu} is the selected accusative plural “acorns.” Do as you wish is not a statement that they have already acted.
% mdparagraph:P130:end

\end{minipage}\par\addvspace{4pt}

\clearpage
% mdparagraph:P131:begin
\section*{Cook 45}
% mdparagraph:P131:end

\par\begin{minipage}{\linewidth}\raggedright
% mdparagraph:P132:begin
\subsection*{Cook 45.1}
% mdparagraph:P132:end

% key: KEY-EX-XII-XIII-45-01
% mdparagraph:P133:begin
Nimað ēowerne brōðor and gāð tō ðǣm męnn.
% mdparagraph:P133:end

% output: EX-XII-XIII-45-01-O001
% mdparagraph:P134:begin
\textbf{Cook’s prompt:} Take your (plur.) brother and go to the man.
% mdparagraph:P134:end

% mdparagraph:P135:begin
\textbf{Target and explanation:} \textbf{Nimað} and \textbf{gāð} address the brothers collectively. The brother taken is singular accusative \textbf{ēowerne brōðor}, with the plural possessors expressed by ēowerne. \textbf{Tō ðǣm męnn} gives their destination. A singular command changes the audience; pluralizing brōðor changes the person they are to take.
% mdparagraph:P135:end

\end{minipage}\par\addvspace{4pt}

\par\begin{minipage}{\linewidth}\raggedright
% mdparagraph:P136:begin
\subsection*{Cook 45.2}
% mdparagraph:P136:end

% key: KEY-EX-XII-XIII-45-02
% mdparagraph:P137:begin
And mīn Dryhten gedō hine glædne wið ēow, ðæt hē ēow ēowerne brōðor āgiefe, ðe hē mid him hæfð; ic eom nū mīnra bearna berēafod.
% mdparagraph:P137:end

% output: EX-XII-XIII-45-02-O001
% mdparagraph:P138:begin
\textbf{Cook’s prompt:} And my Lord make (gedōn) him gracious (glæd) toward (wið) you, that he may give (āgiefan, 196. f) you your brother, whom he hath with him; I am now bereft [despoiled] of my children (156).
% mdparagraph:P138:end

% mdparagraph:P139:begin
\textbf{Target and explanation:} \textbf{Gedō} is a third-singular subjunctive wish, “may my Lord make him …,” not an imperative addressed to the Lord. It has accusative \textbf{hine} and agreeing predicate \textbf{glædne}, here gracious toward them. Purpose \textbf{āgiefe} restores the accusative brother to dative \textbf{ēow}; the relative says that the chief presently has him. \textbf{Eom berēafod} is passive, with genitive \textbf{mīnra bearna}, what Jacob is deprived of. Changing that genitive to a nominative subject would reverse the deprivation.
% mdparagraph:P139:end

\end{minipage}\par\addvspace{4pt}

\par\begin{minipage}{\linewidth}\raggedright
% mdparagraph:P140:begin
\subsection*{Cook 45.3}
% mdparagraph:P140:end

% key: KEY-EX-XII-XIII-45-03
% mdparagraph:P141:begin
Ðā nōmon hīe ðā lāc and twīfeald feoh, and fōron tō Egypta lande tō Iosepe.
% mdparagraph:P141:end

% output: EX-XII-XIII-45-03-O001
% mdparagraph:P142:begin
\textbf{Cook’s prompt:} Then they took the presents and double (twīfeald) money (feoh, 47), and went to the land of the Egyptians to Joseph (Iosep).
% mdparagraph:P142:end

% mdparagraph:P143:begin
\textbf{Target and explanation:} The brothers took both the presents and double money, then travelled to Joseph. \textbf{Nōmon} and \textbf{fōron} are third-plural preterites. \textbf{Twīfeald feoh} is accusative mass singular, not a plural count of coins; retain the source’s ī and unmarked feoh. The whole supplied destination \textbf{tō Egypta lande tō Iosepe} distinguishes the land from Joseph as the person approached. It does not require an independently invented foreign-name declension.
% mdparagraph:P143:end

\end{minipage}\par\addvspace{4pt}

\par\begin{minipage}{\linewidth}\raggedright
% mdparagraph:P144:begin
\subsection*{Cook 45.4}
% mdparagraph:P144:end

% key: KEY-EX-XII-XIII-45-04
% mdparagraph:P145:begin
Ðā hē hīe geseah, ðā cwæð hē tō his gerēfan: Lǣd ðās męnn in, and gearca ūrne gebēorscipe, ðæt hīe mægen mid mē etan tō middes dæges.
% mdparagraph:P145:end

% output: EX-XII-XIII-45-04-O001
% mdparagraph:P146:begin
\textbf{Cook’s prompt:} When he saw them, then said he to (tō) his steward (gerēfa, 53), Lead in these men, and prepare our feast, that they may (mugan) eat with me at noon (tō middes dæges, 158; 166. 1).
% mdparagraph:P146:end

% mdparagraph:P147:begin
\textbf{Target and explanation:} Joseph sees the men and speaks to one steward: \textbf{lǣd} and \textbf{gearca} are singular commands. In the purpose clause \textbf{hīe mægen etan}, the brothers become the third-plural subject; \textbf{mid mē} refers to Joseph. The supplied time phrase \textbf{tō middes dæges} means at noon, with genitive middes agreeing with dæges. The ability form is not a new finite “eat” form, and the noon phrase is not a command to use the ordinary dative after every tō.
% mdparagraph:P147:end

\end{minipage}\par\addvspace{4pt}

\par\begin{minipage}{\linewidth}\raggedright
% mdparagraph:P148:begin
\subsection*{Cook 45.5}
% mdparagraph:P148:end

% key: KEY-EX-XII-XIII-45-05
% mdparagraph:P149:begin
And hē dyde swā him boden wæs.
% mdparagraph:P149:end

% output: EX-XII-XIII-45-05-O001
% mdparagraph:P150:begin
\textbf{Cook’s prompt:} And he did as he was commanded [was commanded him; cf. p. 177, l. 10].
% mdparagraph:P150:end

% mdparagraph:P151:begin
\textbf{Target and explanation:} \textbf{Hē dyde} is the steward’s action. \textbf{Swā him boden wæs} is the supplied impersonal command passive: “as was commanded him.” \textbf{Him} remains dative addressee, not a nominative passive subject. Cook’s cited parallel has hiere beboden wæs; the model uses the separately attested command sense of bēodan. A personal passive paraphrase may be possible, but this exercise expressly supplies the impersonal target.
% mdparagraph:P151:end

\end{minipage}\par\addvspace{4pt}

\par\begin{minipage}{\linewidth}\raggedright
% mdparagraph:P152:begin
\subsection*{Cook 45.6}
% mdparagraph:P152:end

% key: KEY-EX-XII-XIII-45-06
% mdparagraph:P153:begin
Ðā hīe ðǣr wǣron, ðā wurdon hīe swīðe āfǣred, and sægdon betwēonan him: For ðǣm fēo ðe wē on ūrum saccum fundon, wē sindon hider in lǣded, ðæt hē ūs ūrra assena and ūrra ðinga berēafie, and ūs selfe ðēowie.
% mdparagraph:P153:end

% output: EX-XII-XIII-45-06-O001
% mdparagraph:P154:begin
\textbf{Cook’s prompt:} When they were there (ðǣr), then they became greatly frightened, and said to one another (betwēonan him), Because of (for) the money that we found in our sacks (sacc, 43), we (om. have) been led (62) in hither, that he may despoil us of our asses, and of our things, and enslave us (our)selves.
% mdparagraph:P154:end

% mdparagraph:P155:begin
\textbf{Target and explanation:} The narrative’s plural \textbf{wurdon … āfǣred} describes the brothers becoming frightened; a predicative participle may remain uninflected. Inside their speech, \textbf{for ðǣm fēo} is causal, and \textbf{on ūrum saccum} is static location. \textbf{Wē sindon … lǣded} preserves the passive without adding an English-style have; base lǣded is licensed with plural wē. \textbf{Berēafie} takes accusative ūs and genitive possessions \textbf{ūrra assena /\allowbreak{} ūrra ðinga}. \textbf{Ūs selfe} is the object of \textbf{ðēowie}, here enslave, not serve. Reversing those roles reverses their fear.
% mdparagraph:P155:end

\end{minipage}\par\addvspace{4pt}

\par\begin{minipage}{\linewidth}\raggedright
% mdparagraph:P156:begin
\subsection*{Cook 45.7}
% mdparagraph:P156:end

% key: KEY-EX-XII-XIII-45-07
% mdparagraph:P157:begin
Ðā cwǣdon hīe tō ðǣm gerēfan: Lēof, wē biddað ðē ðæt ðū ūrre sprǣce hlyste; wē wǣron hēr hwīlum ǣr and bohton ūs hwǣte.
% mdparagraph:P157:end

% output: EX-XII-XIII-45-07-O001
% mdparagraph:P158:begin
\textbf{Cook’s prompt:} Then said they to (tō) the steward, O sir, we pray (biddan) thee that thou hear (hlystan, 156) our speech; we were here a while ago and bought wheat (hwǣte, 44) for ourselves (161; 184. a).
% mdparagraph:P158:end

% mdparagraph:P159:begin
\textbf{Target and explanation:} \textbf{Lēof} addresses the steward. \textbf{Biddað} is an assertion with wē, “we beseech,” although the spelling can also occur in a command. \textbf{Ðē} is the person beseeched; \textbf{ðū} is subject of subjunctive \textbf{hlyste}, which governs genitive \textbf{ūrre sprǣce}. \textbf{Bohton ūs hwǣte} means “we bought wheat for ourselves”: ūs is benefactive/\allowbreak{}reflexive dative, wheat the accusative object. The source gives no instruction to omit “for ourselves.” Confusing the two pronouns’ roles or making the speech accusative misses the assigned government.
% mdparagraph:P159:end

% mdparagraph:P160:begin
\textbf{Supplied construction:} \textbf{Lēof} means “sir” in this supplied respectful address, and \textbf{hwīlum ǣr} means “a while ago.” These bounded uses are given beside the task; they are not tests of the earlier senses “dear” and “sometimes.” A grammatical alternative preserving respectful address and earlier time may also satisfy Cook’s prompt.
% mdparagraph:P160:end

\end{minipage}\par\addvspace{4pt}

\par\begin{minipage}{\linewidth}\raggedright
% mdparagraph:P161:begin
\subsection*{Cook 45.8}
% mdparagraph:P161:end

% key: KEY-EX-XII-XIII-45-08
% mdparagraph:P162:begin
Ðā wē cēapodon and hāmweard wǣron, ðā undydon wē ūre saccas; ðā fundon wē ðæt feoh ðǣron ðe wē ǣr sealdon.
% mdparagraph:P162:end

% output: EX-XII-XIII-45-08-O001
% mdparagraph:P163:begin
\textbf{Cook’s prompt:} When we had trafficked (cēapian) and were on the way home (hāmweard wǣron), then we undid (undōn) our sacks; then we found the money therein which we had given (see ǣr, adv.).
% mdparagraph:P163:end

% mdparagraph:P164:begin
\textbf{Target and explanation:} The first \textbf{cēapodon} and later \textbf{sealdon} are simple preterites with the anterior meanings “had traded” and “had given”; \textbf{ǣr} makes the earlier payment explicit. \textbf{Hāmweard wǣron} is the supplied homeward frame. \textbf{Undydon} says that the speakers opened their sacks; \textbf{fundon} then states their discovery. \textbf{Ðe} refers to the money previously given, not the sacks. Adding a perfect auxiliary is unnecessary, while removing the earlier-payment relation loses part of the meaning.
% mdparagraph:P164:end

\end{minipage}\par\addvspace{4pt}

\par\begin{minipage}{\linewidth}\raggedright
% mdparagraph:P165:begin
\subsection*{Cook 45.9}
% mdparagraph:P165:end

% key: KEY-EX-XII-XIII-45-09
% mdparagraph:P166:begin
Nū wē hit ęft brōhton, and ēac ōðer seolfor mid tō cēapianne; and wē ne witon hwā hit on ūre saccas dyde.
% mdparagraph:P166:end

% output: EX-XII-XIII-45-09-O001
% mdparagraph:P167:begin
\textbf{Cook’s prompt:} Now we have brought it back, and also other silver to traffic with [with to traffic]; and we know not who put (dōn) it into (on, 172) our sacks.
% mdparagraph:P167:end

% mdparagraph:P168:begin
\textbf{Target and explanation:} \textbf{Brōhton} expresses the bringing back without requiring an auxiliary have. “Other silver” is a second object of that bringing; \textbf{mid tō cēapianne} is the supplied compressed instrument-and-purpose phrase, retaining the requested “with to trade” order. \textbf{Hwā} introduces the indirect question about the unknown person who \textbf{dyde} the money into the sacks. The movement/\allowbreak{}destination phrase \textbf{on ūre saccas} is accusative plural. Do not interpret the silver as a human trading partner or the final question as a new direct question.
% mdparagraph:P168:end

\end{minipage}\par\addvspace{4pt}

\clearpage
% mdparagraph:P169:begin
\section*{Cook 46}
% mdparagraph:P169:end

\par\begin{minipage}{\linewidth}\raggedright
% mdparagraph:P170:begin
\subsection*{Cook 46.1}
% mdparagraph:P170:end

% key: KEY-EX-XII-XIII-46-01
% mdparagraph:P171:begin
Ðā cwæð se gerēfa: Sibb sīe mid ēow; ne ondrǣdað ēow; ēower God and ēowres fæder God sealde ēow goldhord on ēowrum saccum.
% mdparagraph:P171:end

% output: EX-XII-XIII-46-01-O001
% mdparagraph:P172:begin
\textbf{Cook’s prompt:} Then said the steward, Peace be with you, fear (ondrǣdan, 184) ye not; your God and your father’s God (om. hath) given you treasure in your sacks.
% mdparagraph:P172:end

% mdparagraph:P173:begin
\textbf{Target and explanation:} \textbf{Sibb sīe mid ēow} is a wish of peace. \textbf{Ondrǣdað ēow} is a negative plural command with reflexive dative, no expressed feared object. The two descriptions “your God” and “your father’s God” refer here to one deity, hence singular \textbf{sealde}. \textbf{Ēowres} agrees with genitive fæder; the following \textbf{ēow} is dative recipient of the treasure. The money is located \textbf{on ēowrum saccum}, dative plural. Making two gods or changing the possessive’s case changes the supplied interpretation.
% mdparagraph:P173:end

\end{minipage}\par\addvspace{4pt}

\par\begin{minipage}{\linewidth}\raggedright
% mdparagraph:P174:begin
\subsection*{Cook 46.2}
% mdparagraph:P174:end

% key: KEY-EX-XII-XIII-46-02
% mdparagraph:P175:begin
And hīe ðwōgon hiera fēt, and hē sealde hiera assum fōdor.
% mdparagraph:P175:end

% output: EX-XII-XIII-46-02-O001
% mdparagraph:P176:begin
\textbf{Cook’s prompt:} And they washed (101) their feet, and he gave their asses provender [fodder].
% mdparagraph:P176:end

% mdparagraph:P177:begin
\textbf{Target and explanation:} \textbf{Ðwōgon} is the contracted verb’s third-plural preterite: the brothers washed their accusative plural \textbf{fēt}. The subject then changes to singular \textbf{hē}, the steward, who \textbf{sealde} fodder to \textbf{hiera assum}, dative plural. \textbf{Hiera} is a possessor in both phrases; its unchanged spelling does not give the feet and asses the same grammatical role.
% mdparagraph:P177:end

\end{minipage}\par\addvspace{4pt}

\par\begin{minipage}{\linewidth}\raggedright
% mdparagraph:P178:begin
\subsection*{Cook 46.3}
% mdparagraph:P178:end

% key: KEY-EX-XII-XIII-46-03
% mdparagraph:P179:begin
Hīe lęgdon hiera lāc forð, ongēan ðætte Iosep inēode.
% mdparagraph:P179:end

% output: EX-XII-XIII-46-03-O001
% mdparagraph:P180:begin
\textbf{Cook’s prompt:} They set (lęcgean) forth their presents, against that (ongēan ðætte) Joseph (om. should) enter.
% mdparagraph:P180:end

% mdparagraph:P181:begin
\textbf{Target and explanation:} \textbf{Lęgdon} is third-plural preterite of lęcgean, with accusative plural \textbf{lāc} and forð “forth.” The supplied \textbf{ongēan ðætte} clause means in readiness for Joseph’s entering. \textbf{Inēode} is past subjunctive third singular in that prospective frame; its identical indicative spelling elsewhere does not settle its mood here. A mere assertion that he had already entered would lose the preparation relation.
% mdparagraph:P181:end

\end{minipage}\par\addvspace{4pt}

\par\begin{minipage}{\linewidth}\raggedright
% mdparagraph:P182:begin
\subsection*{Cook 46.4}
% mdparagraph:P182:end

% key: KEY-EX-XII-XIII-46-04
% mdparagraph:P183:begin
Ðā oncnēow Iosep hīe, and āscode hīe hwæðer hiera fæder hāl wǣre, oððe hwæðer hē lifde.
% mdparagraph:P183:end

% output: EX-XII-XIII-46-04-O001
% mdparagraph:P184:begin
\textbf{Cook’s prompt:} Then Joseph recognized them, and asked them whether their father was hale, or whether he was living [lived].
% mdparagraph:P184:end

% mdparagraph:P185:begin
\textbf{Target and explanation:} Joseph is the singular recognizer and questioner; both accusative \textbf{hīe} forms refer to the brothers. \textbf{Oncnēow} is the preterite of recognition. In the indirect questions the subject is their father, with third-singular past subjunctives \textbf{wǣre} and \textbf{lifde}. The question asks whether he was well or alive; it does not assert that he was dead or make the brothers the subject of lifde.
% mdparagraph:P185:end

\end{minipage}\par\addvspace{4pt}

\par\begin{minipage}{\linewidth}\raggedright
% mdparagraph:P186:begin
\subsection*{Cook 46.5}
% mdparagraph:P186:end

% key: KEY-EX-XII-XIII-46-05
% mdparagraph:P187:begin
Ðā cwǣdon hīe: Ðīn ðēow, ūre fæder, is gesund; hē leofað gīet.
% mdparagraph:P187:end

% output: EX-XII-XIII-46-05-O001
% mdparagraph:P188:begin
\textbf{Cook’s prompt:} Then said they, Thy servant (ðēow) our father is well (gesund); he is yet alive [liveth yet].
% mdparagraph:P188:end

% mdparagraph:P189:begin
\textbf{Target and explanation:} “Thy servant, our father” names one person, the singular subject of \textbf{is} and \textbf{leofað}. \textbf{Gesund} means well or healthy; \textbf{gīet} preserves “yet.” \textbf{Ðēow} is the selected strong servant noun, not the weak noun ðēowa or an adjective silently substituted for it. A present assertion of life is needed; past lifde would answer a different temporal question.
% mdparagraph:P189:end

\end{minipage}\par\addvspace{4pt}

\par\begin{minipage}{\linewidth}\raggedright
% mdparagraph:P190:begin
\subsection*{Cook 46.6}
% mdparagraph:P190:end

% key: KEY-EX-XII-XIII-46-06
% mdparagraph:P191:begin
Ðā Iosep geseah his brōðor Beniamin, ðā cwæð hē: Is ðis se cnapa ðe gē mē foresægdon? And hē cwæð: God gemiltsie ðē, mīn sunu.
% mdparagraph:P191:end

% output: EX-XII-XIII-46-06-O001
% mdparagraph:P192:begin
\textbf{Cook’s prompt:} When Joseph saw his brother Benjamin, then said he, Is this (neut.) the lad, (om. of) whom (87. d) ye told me before (foresęcgean)? And he said, God have compassion on thee, my son.
% mdparagraph:P192:end

% mdparagraph:P193:begin
\textbf{Target and explanation:} Joseph first speaks of Benjamin in the identification question, then addresses him directly. Retain the requested neuter \textbf{ðis} in \textbf{is ðis se cnapa}. The relative \textbf{ðe} represents the boy mentioned earlier; \textbf{mē} is dative hearer after \textbf{foresægdon}, and no “of” is inserted. \textbf{Gemiltsie} is a third-singular subjunctive wish with dative \textbf{ðē}, “may God have compassion on you.” \textbf{Mīn sunu} is address, not the subject of that wish.
% mdparagraph:P193:end

\end{minipage}\par\addvspace{4pt}

\par\begin{minipage}{\linewidth}\raggedright
% mdparagraph:P194:begin
\subsection*{Cook 46.7}
% mdparagraph:P194:end

% key: KEY-EX-XII-XIII-46-07
% mdparagraph:P195:begin
Ðā bēad Iosep his gerēfan, and cwæð: Gefyl hiera saccas mid hwǣte, and sęte ǣlces hiera feoh on his āgenne sacc.
% mdparagraph:P195:end

% output: EX-XII-XIII-46-07-O001
% mdparagraph:P196:begin
\textbf{Cook’s prompt:} Then Joseph commanded (bēodan, 164) his steward, and said, Fill their sacks with wheat, and put the money of each of them into (on, 172) his own sack.
% mdparagraph:P196:end

% mdparagraph:P197:begin
\textbf{Target and explanation:} Commanding \textbf{bēad} takes dative \textbf{his gerēfan}. Singular \textbf{gefyl} and \textbf{sęte} address that steward. \textbf{Mid hwǣte} supplies the sacks’ contents, while \textbf{ǣlces hiera feoh} means each one’s money: the money belongs to each brother. \textbf{On his āgenne sacc} is the accusative destination, each owner’s own sack. Making all the money belong to the steward or using a static dative for the placement changes the instructions.
% mdparagraph:P197:end

% mdparagraph:P198:begin
\textbf{Supplied construction:} The visible frame \textbf{Gefyl hiera saccas mid hwǣte} supplies the material construction: \textbf{saccas} is accusative plural and \textbf{hwǣte} dative singular after \textbf{mid}. Do not grade recall of this previously untaught lexical frame.
% mdparagraph:P198:end

\end{minipage}\par\addvspace{4pt}

\par\begin{minipage}{\linewidth}\raggedright
% mdparagraph:P199:begin
\subsection*{Cook 46.8}
% mdparagraph:P199:end

% key: KEY-EX-XII-XIII-46-08
% mdparagraph:P200:begin
And nim mīnne seolfrenne læfel, and ðæs hwǣtes weorð ðe hē ðē sealde, and dō on ðæs gingstan sacc; and hē dyde swā.
% mdparagraph:P200:end

% output: EX-XII-XIII-46-08-O001
% mdparagraph:P201:begin
\textbf{Cook’s prompt:} And take my silver cup (læfel, 43), and the price (weorð, 47) of the wheat, which he gave thee, and put into the sack of the youngest; and he did so.
% mdparagraph:P201:end

% mdparagraph:P202:begin
\textbf{Target and explanation:} The two things to take and put are the silver cup and the wheat-price. Supplied \textbf{seolfrenne} agrees with masculine accusative singular \textbf{læfel}. \textbf{Ðæs hwǣtes} specifies the price, and the relative says that he gave it \textbf{ðē}, to thee: preserve the source’s singular addressee. The second imperative \textbf{dō} shares the coordinated objects without an inserted singular “it.” \textbf{Ðæs gingstan} is the substantival weak genitive, the youngest person as possessor of the sack.
% mdparagraph:P202:end

\end{minipage}\par\addvspace{4pt}

\par\begin{minipage}{\linewidth}\raggedright
% mdparagraph:P203:begin
\subsection*{Cook 46.9}
% mdparagraph:P203:end

% key: KEY-EX-XII-XIII-46-09
% mdparagraph:P204:begin
On morgen, ðā hīe fōron, and hīe būtan byrig wǣron, and fōron sumne dǣl ðæs weges, ðā cwæð Iosep tō his gerēfan: Ārīs, and gā æfter ðissum mannum; and ðonne ðū hīe fō, ðonne āsca hīe hwȳ hīe wolden gōd mid yfele forgieldan.
% mdparagraph:P204:end

% output: EX-XII-XIII-46-09-O001
% mdparagraph:P205:begin
\textbf{Cook’s prompt:} In (om. the) morning, as (ðā) they journeyed, and they were without (om. the) city, and had gone (faran) some part of the way, then said Joseph to (tō) his steward, Arise and go after these men, and when you have (196. b) overtaken [caught] them, then ask them why they would (willan) requite good with (mid) evil.
% mdparagraph:P205:end

% mdparagraph:P206:begin
\textbf{Target and explanation:} Retain all three circumstances: they travelled, were outside the city, and had gone part of the way. \textbf{Būtan byrig} is the supplied spatial outside phrase; \textbf{sumne dǣl ðæs weges} expresses distance travelled, and the simple preterite fōron can convey “had gone.” \textbf{Ārīs, gā, āsca} address one steward. Future temporal \textbf{ðonne ðū hīe fō} has second-singular present subjunctive, “when you have caught them.” In the indirect why-question, \textbf{wolden} is plural past subjunctive and \textbf{gōd mid yfele forgieldan} means repay good with evil. The evil is what is given in return.
% mdparagraph:P206:end

% mdparagraph:P207:begin
\textbf{Supplied construction:} The fixed phrase \textbf{On morgen} is supplied with accusative singular \textbf{morgen}. Its temporal case is part of the visible help, not an inference from the earlier numbered-day construction.
% mdparagraph:P207:end

\end{minipage}\par\addvspace{4pt}

\clearpage
% mdparagraph:P208:begin
\section*{Cook 47}
% mdparagraph:P208:end

\par\begin{minipage}{\linewidth}\raggedright
% mdparagraph:P209:begin
\subsection*{Cook 47.1}
% mdparagraph:P209:end

% key: KEY-EX-XII-XIII-47-01
% mdparagraph:P210:begin
Se læfel ðe gē forstǣlon wæs mīnum hlāforde swīðe dēorwierðe.
% mdparagraph:P210:end

% output: EX-XII-XIII-47-01-O001
% mdparagraph:P211:begin
\textbf{Cook’s prompt:} The cup, which ye stole (forstelan, 105), was very precious to my lord (hlāford).
% mdparagraph:P211:end

% mdparagraph:P212:begin
\textbf{Target and explanation:} The cup is the singular subject of \textbf{wæs}. Relative \textbf{ðe} represents it as object of \textbf{forstǣlon}, past indicative second plural addressed to the brothers. \textbf{Mīnum hlāforde} is dative, the person to whom it was precious. The past plural’s \textbf{ǣ} differs from singular forstæl in 47.8. The small isolated mark in the print does not supply an additional word or grammatical instruction.
% mdparagraph:P212:end

\end{minipage}\par\addvspace{4pt}

\par\begin{minipage}{\linewidth}\raggedright
% mdparagraph:P213:begin
\subsection*{Cook 47.2}
% mdparagraph:P213:end

% key: KEY-EX-XII-XIII-47-02
% mdparagraph:P214:begin
Hē dyde swā him boden wæs, and ðā hīe fangen wǣron, ðā cwǣdon hīe: Ðæt feoh ðe wē on ūrum saccum fundon, wē brōhton ðē of ūrum lande; wēnst ðū ðæt wē ðīnes hlāfordes gold oððe his seolfor stǣlen?
% mdparagraph:P214:end

% output: EX-XII-XIII-47-02-O001
% mdparagraph:P215:begin
\textbf{Cook’s prompt:} He did as was commanded him, and when they were overtaken [caught], they said, The money that we found in our sacks we brought to thee from our country; thinkest (wēnan) thou that we (om. have) stolen (stelan) thy lord’s gold or his silver?
% mdparagraph:P215:end

% mdparagraph:P216:begin
\textbf{Target and explanation:} The steward did what was commanded him, using the same impersonal frame as 45.5. \textbf{Fangen wǣron} is passive with the brothers as plural subject. Their following claim preserves the money, its discovery in the sacks, the return to the steward and the origin in their country. \textbf{Wēnst ðū ðæt … stǣlen} asks whether he thinks they stole the gold or silver. \textbf{Stǣlen} is first-plural past subjunctive; it supplies the questioned earlier theft without adding the omitted have. It is not an admission that they stole anything.
% mdparagraph:P216:end

\end{minipage}\par\addvspace{4pt}

\par\begin{minipage}{\linewidth}\raggedright
% mdparagraph:P217:begin
\subsection*{Cook 47.3}
% mdparagraph:P217:end

% key: KEY-EX-XII-XIII-47-03
% mdparagraph:P218:begin
Sēc hit; and sīe hē ofslægen mid ðǣm ðe ðū hit finde.
% mdparagraph:P218:end

% output: EX-XII-XIII-47-03-O001
% mdparagraph:P219:begin
\textbf{Cook’s prompt:} Look for [seek out] it; and let him be slain (ofslēan) with whom thou findest (195) it.
% mdparagraph:P219:end

% mdparagraph:P220:begin
\textbf{Target and explanation:} \textbf{Sēc} is a singular command. \textbf{Sīe hē ofslægen} is a third-singular wish, “let him be slain”; it is not an instruction that the addressee kill himself. In \textbf{mid ðǣm ðe ðū hit finde}, mid governs the dative associate and invariant ðe resumes that person. \textbf{Finde} is second-singular present subjunctive in the optative/\allowbreak{}hypothetical relative under Cook 195. \textbf{Hit} is neuter accusative singular here, resuming the alleged stolen gold or silver in 47.2 (both neuter). A version directly resuming masculine \textbf{læfel} instead uses \textbf{hine}; preserve its reference consistently.
% mdparagraph:P220:end

\end{minipage}\par\addvspace{4pt}

\par\begin{minipage}{\linewidth}\raggedright
% mdparagraph:P221:begin
\subsection*{Cook 47.4}
% mdparagraph:P221:end

% key: KEY-EX-XII-XIII-47-04
% mdparagraph:P222:begin
Ðā cwæð hē: Mid swā hwǣm swā ic hit finde, sīe hē mīn ðēow, and sīen ðā ōðre clǣne.
% mdparagraph:P222:end

% output: EX-XII-XIII-47-04-O001
% mdparagraph:P223:begin
\textbf{Cook’s prompt:} Then said he, With whomsoever I find it, let him be my servant, and let the others be blameless (clǣne).
% mdparagraph:P223:end

% mdparagraph:P224:begin
\textbf{Target and explanation:} The free-choice phrase \textbf{mid swā hwǣm swā} identifies whichever person has the alleged stolen property; neuter \textbf{hit} continues 47.2’s gold/\allowbreak{}silver reference, while direct reference to \textbf{læfel} would use \textbf{hine}. \textbf{Hē} resumes that person as subject of singular \textbf{sīe}. The remaining brothers are plural \textbf{ðā ōðre} with \textbf{sīen}, and \textbf{clǣne} means clear of blame here. “Whoever I find it with” is an equivalent English rendering; making all the brothers guilty or making the indefinite person plural reverses the contrast.
% mdparagraph:P224:end

\end{minipage}\par\addvspace{4pt}

\par\begin{minipage}{\linewidth}\raggedright
% mdparagraph:P225:begin
\subsection*{Cook 47.5}
% mdparagraph:P225:end

% key: KEY-EX-XII-XIII-47-05
% mdparagraph:P226:begin
And hē sōhte fram ðǣm ieldestan oð ðone gingstan; and hē funde ðone læfel on Beniamines sacce; and hīe ęft gewęndon tō ðǣre byrig.
% mdparagraph:P226:end

% output: EX-XII-XIII-47-05-O001
% mdparagraph:P227:begin
\textbf{Cook’s prompt:} And he sought from (fram) the eldest unto [until] the youngest; and he found the cup in Benjamin’s sack; and they turned back to the city.
% mdparagraph:P227:end

% mdparagraph:P228:begin
\textbf{Target and explanation:} The search proceeds \textbf{fram ðǣm ieldestan oð ðone gingstan}: dative after fram, accusative after oð. \textbf{Funde} is the already taught and source-attested third-singular form of findan; it is not an error or necessarily second singular. \textbf{Fand} is also attested, but funde is the selected model. The cup is found in supplied \textbf{Beniamines sacce}, static dative; plural \textbf{gewęndon} then returns the brothers to the city. Do not make the steward the plural returning subject.
% mdparagraph:P228:end

\end{minipage}\par\addvspace{4pt}

\par\begin{minipage}{\linewidth}\raggedright
% mdparagraph:P229:begin
\subsection*{Cook 47.6}
% mdparagraph:P229:end

% key: KEY-EX-XII-XIII-47-06
% mdparagraph:P230:begin
And Iudas inēode fyrmest mid his brōðrum tō Iosepe, and hīe ealle ætgædere beforan Iosepe fēollon.
% mdparagraph:P230:end

% output: EX-XII-XIII-47-06-O001
% mdparagraph:P231:begin
\textbf{Cook’s prompt:} And Judah went in first (fyrmest) with his brethren to Joseph, and they fell all together before Joseph.
% mdparagraph:P231:end

% mdparagraph:P232:begin
\textbf{Target and explanation:} \textbf{Iudas inēode fyrmest} says that Judah entered first, with his brothers. \textbf{Mid his brōðrum} describes his companions, not a plural subject of inēode. The next clause expressly has \textbf{hīe ealle} and plural \textbf{fēollon}: all fell together before Joseph. Both the leading entry and the collective prostration must be retained.
% mdparagraph:P232:end

\end{minipage}\par\addvspace{4pt}

\par\begin{minipage}{\linewidth}\raggedright
% mdparagraph:P233:begin
\subsection*{Cook 47.7}
% mdparagraph:P233:end

% key: KEY-EX-XII-XIII-47-07
% mdparagraph:P234:begin
And hē cwæð tō him: Hwȳ gē swā dōn woldon?
% mdparagraph:P234:end

% output: EX-XII-XIII-47-07-O001
% mdparagraph:P235:begin
\textbf{Cook’s prompt:} And he said to (tō) them, Why would ye do so (verb last)?
% mdparagraph:P235:end

% mdparagraph:P236:begin
\textbf{Target and explanation:} The direct question keeps \textbf{gē} as the plural addressees and \textbf{woldon} as the final finite verb, governing plain \textbf{dōn}. \textbf{Swā} means “so/\allowbreak{}in that manner,” referring to the alleged conduct. A grammatical question with a different order can miss the explicit verb-last control; a singular modal would change the audience.
% mdparagraph:P236:end

\end{minipage}\par\addvspace{4pt}

\par\begin{minipage}{\linewidth}\raggedright
% mdparagraph:P237:begin
\subsection*{Cook 47.8}
% mdparagraph:P237:end

% key: KEY-EX-XII-XIII-47-08
% mdparagraph:P238:begin
Sīe hē mīn ðēow ðe ðone læfel forstæl, and gāð gē frige tō ēowrum fæder.
% mdparagraph:P238:end

% output: EX-XII-XIII-47-08-O001
% mdparagraph:P239:begin
\textbf{Cook’s prompt:} Let him be my servant who stole the cup, and go ye free to your father.
% mdparagraph:P239:end

% mdparagraph:P240:begin
\textbf{Target and explanation:} The singular wish \textbf{sīe hē mīn ðēow} is restricted by the relative to the one who stole the cup. \textbf{Forstæl} is past indicative third singular, unlike plural forstǣlon in 47.1. The rest receive plural command \textbf{gāð}, with supplied \textbf{frige}, free. Do not change the selected source spelling by silently adding vowel length, and do not extend the servitude to those being released.
% mdparagraph:P240:end

\end{minipage}\par\addvspace{4pt}

\par\begin{minipage}{\linewidth}\raggedright
% mdparagraph:P241:begin
\subsection*{Cook 47.9}
% mdparagraph:P241:end

% key: KEY-EX-XII-XIII-47-09
% mdparagraph:P242:begin
Ðā cwæð Iudas: Ðū āscodest ūs ǣr hwæðer wē fæder oððe brōðor hæfden.
% mdparagraph:P242:end

% output: EX-XII-XIII-47-09-O001
% mdparagraph:P243:begin
\textbf{Cook’s prompt:} Then said Judah, Thou askedst us before whether we had father or brother.
% mdparagraph:P243:end

% mdparagraph:P244:begin
\textbf{Target and explanation:} Judah addresses Joseph with singular \textbf{ðū āscodest}. \textbf{Ūs} is the plural people questioned; \textbf{hwæðer wē … hæfden} is the embedded question about their family, with first-plural past subjunctive. \textbf{Ǣr} retains the earlier time. Converting the embedded question into a present assertion, or making Jacob the person who asked, changes the report.
% mdparagraph:P244:end

\end{minipage}\par\addvspace{4pt}

\par\begin{minipage}{\linewidth}\raggedright
% mdparagraph:P245:begin
\subsection*{Cook 47.10}
% mdparagraph:P245:end

% key: KEY-EX-XII-XIII-47-10
% mdparagraph:P246:begin
And wē cwǣdon: Wē habbað ealdne fæder, and hē hæfð mid him ūrne gingstan brōðor; and hē lufað hine ænne ofer ūs ealle.
% mdparagraph:P246:end

% output: EX-XII-XIII-47-10-O001
% mdparagraph:P247:begin
\textbf{Cook’s prompt:} And we said, We have (om. an) old father, and he hath with him our youngest brother, and he loveth him alone (ænne) above (ofer) us all.
% mdparagraph:P247:end

% mdparagraph:P248:begin
\textbf{Target and explanation:} The brothers are the subject of \textbf{habbað}, with accusative \textbf{ealdne fæder} and no indefinite article. The father then becomes singular subject of \textbf{hæfð} and \textbf{lufað}. \textbf{Ænne}, as supplied, agrees with accusative hine in the sense “him alone”; it is not the article omitted before father. \textbf{Ofer ūs ealle} preserves the comparison “above all of us.” Do not replace the singular youngest brother with several brothers.
% mdparagraph:P248:end

\end{minipage}\par\addvspace{4pt}

\par\begin{minipage}{\linewidth}\raggedright
% mdparagraph:P249:begin
\subsection*{Cook 47.11}
% mdparagraph:P249:end

% key: KEY-EX-XII-XIII-47-11
% mdparagraph:P250:begin
And ðū bīetst ūs ðæt wē hine ðē bringen, ðæt ðū hine gesēo, and wite be ðǣm ðæt wē ne wǣren scēaweras.
% mdparagraph:P250:end

% output: EX-XII-XIII-47-11-O001
% mdparagraph:P251:begin
\textbf{Cook’s prompt:} And thou commandest us that we (om. should) bring him to thee, that thou (om. mightest) see him, and know by (be) that that we were not spies (scēawere, 44).
% mdparagraph:P251:end

% mdparagraph:P252:begin
\textbf{Target and explanation:} Keep the actual present \textbf{bīetst}, “thou commandest.” \textbf{Bringen} is first-plural present subjunctive in the commanded action; \textbf{gesēo} and \textbf{wite} are coordinated second-singular present subjunctives expressing its purpose. \textbf{Be ðǣm} means by that evidence, followed by the content \textbf{ðæt wē ne wǣren scēaweras}, with past subjunctive preserving “were.” The same-looking bīetst of bēatan is not meant. A silent backshift of commandest, or dropping the knowing clause, loses an explicit part of the prompt.
% mdparagraph:P252:end

\end{minipage}\par\addvspace{4pt}

\clearpage
% mdparagraph:P253:begin
\section*{Cook 48}
% mdparagraph:P253:end

\par\begin{minipage}{\linewidth}\raggedright
% mdparagraph:P254:begin
\subsection*{Cook 48.1}
% mdparagraph:P254:end

% key: KEY-EX-XII-XIII-48-01
% mdparagraph:P255:begin
And ic swōr āðas mīnum fæder ðæt ic hine ęft hām tō him brōhte, and cwæð: Būtan ic hine ęft tō ðē bringe, sīe ic ā scyldig wið ðē.
% mdparagraph:P255:end

% output: EX-XII-XIII-48-01-O001
% mdparagraph:P256:begin
\textbf{Cook’s prompt:} And I swore oaths (om. to) my father that I (om. would) bring him home again to (tō) him, and said, Except I bring him again back to thee, let me be ever guilty (scyldig) toward (wið) thee.
% mdparagraph:P256:end

% mdparagraph:P257:begin
\textbf{Target and explanation:} Judah swore accusative \textbf{āðas} to dative \textbf{mīnum fæder}, without tō before that recipient. \textbf{Brōhte} is first-singular past subjunctive in the promised action. In the subsequent direct quotation \textbf{bringe} is first-singular present subjunctive after conditional būtan; \textbf{sīe ic ā scyldig} is the enduring self-imprecation. The father referred to as him in the narrative becomes \textbf{ðē} in the speech addressed to him. The change of person must not create a new recipient.
% mdparagraph:P257:end

\end{minipage}\par\addvspace{4pt}

\par\begin{minipage}{\linewidth}\raggedright
% mdparagraph:P258:begin
\subsection*{Cook 48.2}
% mdparagraph:P258:end

% key: KEY-EX-XII-XIII-48-02
% mdparagraph:P259:begin
Mīn hlāford, lǣt ðone cnapan mid his brōðrum hām gān, and sīe ic ðīn ðēow for hine.
% mdparagraph:P259:end

% output: EX-XII-XIII-48-02-O001
% mdparagraph:P260:begin
\textbf{Cook’s prompt:} My lord, let (lǣtan) the lad go home with his brethren, and (om. let) me be thy servant instead of (for, 172) him.
% mdparagraph:P260:end

% mdparagraph:P261:begin
\textbf{Target and explanation:} Singular \textbf{lǣt} addresses Joseph. \textbf{Ðone cnapan} is accusative actor of plain infinitive \textbf{gān}: permit the boy to go home with his brothers. The second clause is an independent first-singular wish, \textbf{sīe ic ðīn ðēow}, with no repeated lexical “let.” \textbf{For hine} is substitution with accusative, “in his place.” Making Judah the person allowed to depart would reverse his offer.
% mdparagraph:P261:end

\end{minipage}\par\addvspace{4pt}

\par\begin{minipage}{\linewidth}\raggedright
% mdparagraph:P262:begin
\subsection*{Cook 48.3}
% mdparagraph:P262:end

% key: KEY-EX-XII-XIII-48-03
% mdparagraph:P263:begin
Ic ne dearr hām gān būtan ðǣm cnapan, ðȳ-lǣs-ðe ic mīnes fæder sārnesse gesēo.
% mdparagraph:P263:end

% output: EX-XII-XIII-48-03-O001
% mdparagraph:P264:begin
\textbf{Cook’s prompt:} I dare not go home without the lad, lest I see my father’s grief.
% mdparagraph:P264:end

% mdparagraph:P265:begin
\textbf{Target and explanation:} \textbf{Ne dearr … gān} expresses the speaker’s present refusal through lack of daring, with a plain infinitive. \textbf{Būtan ðǣm cnapan} means without the boy, a dative prepositional phrase; it is not the unless conjunction used elsewhere. Negative-purpose \textbf{ðȳ-lǣs-ðe … gesēo} means “lest I see,” with first-singular present subjunctive. The father is genitive possessor of accusative grief. The feared sight is prospective, not an event that has already occurred.
% mdparagraph:P265:end

\end{minipage}\par\addvspace{4pt}

\par\begin{minipage}{\linewidth}\raggedright
% mdparagraph:P266:begin
\subsection*{Cook 48.4}
% mdparagraph:P266:end

% key: KEY-EX-XII-XIII-48-04
% mdparagraph:P267:begin
Ðā ne meahte Iosep hine lęng diernan, ac hē ādrāf ealle ðā ōðre ūt, ðæt nān ęlðēodig mann betweox him ne wǣre.
% mdparagraph:P267:end

% output: EX-XII-XIII-48-04-O001
% mdparagraph:P268:begin
\textbf{Cook’s prompt:} Then could (mugan) not Joseph conceal (diernan, 113) himself longer, but he drove all the others out, that no foreign man (om. should) be among them.
% mdparagraph:P268:end

% mdparagraph:P269:begin
\textbf{Target and explanation:} Past \textbf{meahte} governs plain \textbf{diernan}; accusative \textbf{hine} refers reflexively to Joseph. \textbf{Lęng} modifies the duration of concealment. After \textbf{ādrāf … ūt}, the final clause gives the purpose that no foreign person be among them, with \textbf{wǣre} past subjunctive. \textbf{Nān … ne} is concordant negation. Losing either the expulsion or the purpose reduces the sentence to only part of the event.
% mdparagraph:P269:end

\end{minipage}\par\addvspace{4pt}

\par\begin{minipage}{\linewidth}\raggedright
% mdparagraph:P270:begin
\subsection*{Cook 48.5}
% mdparagraph:P270:end

% key: KEY-EX-XII-XIII-48-05
% mdparagraph:P271:begin
And hē wēop and clēopode hlūdre stefne; and hē cwæð tō his brōðrum: Ic eom Iosep; leofað ūre fæder nū gīet?
% mdparagraph:P271:end

% output: EX-XII-XIII-48-05-O001
% mdparagraph:P272:begin
\textbf{Cook’s prompt:} And he wept and cried (clēopian) with a loud (hlūd) voice (160. 1); and he said to (tō) his brethren, I am Joseph; (om. doth) our father live yet (nū gīet)?
% mdparagraph:P272:end

% mdparagraph:P273:begin
\textbf{Target and explanation:} Joseph both \textbf{wēop} and \textbf{clēopode}. The latter retains the long ē of Cook’s exact clēopian hint; it does not rewrite the previously taught unmarked cleopian spelling. \textbf{Hlūdre stefne} is dative/\allowbreak{}instrumental manner, with a loud voice, without an extra preposition. The quoted \textbf{eom} identifies Joseph; \textbf{leofað … nū gīet?} asks whether the father still lives, without an English-style do auxiliary.
% mdparagraph:P273:end

\end{minipage}\par\addvspace{4pt}

\par\begin{minipage}{\linewidth}\raggedright
% mdparagraph:P274:begin
\subsection*{Cook 48.6}
% mdparagraph:P274:end

% key: KEY-EX-XII-XIII-48-06
% mdparagraph:P275:begin
Ne ondrǣdað ēow nānes ðinges, ne ofðynce ēow ðæt gē mē on ðis rīce sealdon; sōðlice God sęnde mē for ēowerre ðearfe tō Egypta lande.
% mdparagraph:P275:end

% output: EX-XII-XIII-48-06-O001
% mdparagraph:P276:begin
\textbf{Cook’s prompt:} Fear ye (184) nothing, nor (om. let) it grieve (190) you that ye sold (sęllan) me into this kingdom; indeed (om. it was) for your benefit God sent me into the land of the Egyptians.
% mdparagraph:P276:end

% mdparagraph:P277:begin
\textbf{Target and explanation:} \textbf{Ondrǣdað ēow nānes ðinges} combines reflexive dative ēow with genitive of the thing feared. \textbf{Ofðynce ēow} is third-singular impersonal subjunctive, “let it not grieve you”; the plural experiencers do not make its verb plural. \textbf{Sealdon} here means sold, not gave. \textbf{God sęnde} supplies the asserted agent without an added cleft “it was.” The supplied benefit phrase \textbf{for ēowerre ðearfe} is dative feminine singular, including the possessive ending. Replacing it with accusative ēowre would miss that agreement.
% mdparagraph:P277:end

\end{minipage}\par\addvspace{4pt}

\par\begin{minipage}{\linewidth}\raggedright
% mdparagraph:P278:begin
\subsection*{Cook 48.7}
% mdparagraph:P278:end

% key: KEY-EX-XII-XIII-48-07
% mdparagraph:P279:begin
Nū twā gēar hunger wæs ofer ealle eorðan, and gīet sculon fīf, on ðǣm mann nāhwæðer ne ęrian ne rīpan mæg.
% mdparagraph:P279:end

% output: EX-XII-XIII-48-07-O001
% mdparagraph:P280:begin
\textbf{Cook’s prompt:} Now two years has been [was] famine over (172) all (om. the) earth, and (om. there) yet shall (om. be) five, in which one can (mugan) neither (89. a) plow (ęrian, 116) nor reap (rīpan, 102).
% mdparagraph:P280:end

% mdparagraph:P281:begin
\textbf{Target and explanation:} Use the complete supplied frame. \textbf{Twā gēar} gives duration, while singular \textbf{hunger} governs \textbf{wæs}; \textbf{ofer ealle eorðan} is accusative extent without an article. In \textbf{gīet sculon fīf}, fīf stands for five further years and the source explicitly omits “there” and “be.” Do not invent an additional written copula. \textbf{On ðǣm} refers to those years; indefinite \textbf{mann} is singular subject of \textbf{mæg}. \textbf{Nāhwæðer ne ęrian ne rīpan} excludes both activities. “Neither … nor” must not become “either … or.”
% mdparagraph:P281:end

\end{minipage}\par\addvspace{4pt}

\par\begin{minipage}{\linewidth}\raggedright
% mdparagraph:P282:begin
\subsection*{Cook 48.8}
% mdparagraph:P282:end

% key: KEY-EX-XII-XIII-48-08
% mdparagraph:P283:begin
Gāð hrædlīce tō mīnum fæder and sęcgað him ðæt God mē tō hlāforde eallum Egyptum gesętte; bēodað him ðæt hē tō mē cume and on ðissum lande wunie.
% mdparagraph:P283:end

% output: EX-XII-XIII-48-08-O001
% mdparagraph:P284:begin
\textbf{Cook’s prompt:} Go quickly to my father and tell him that God has set me as a (tō) lord (om. to) all (om. the) Egyptians; bid (bēodan) him (164) that he come to me and dwell in this land.
% mdparagraph:P284:end

% mdparagraph:P285:begin
\textbf{Target and explanation:} Plural \textbf{gāð, sęcgað, bēodað} address the brothers. In the supplied appointment clause, \textbf{mē} is the person appointed, \textbf{tō hlāforde} the role, and \textbf{eallum Egyptum} the people governed; the last phrase has neither an added article nor tō. The father is dative \textbf{him} after bēodað, then singular subject of subjunctives \textbf{cume} and \textbf{wunie}. The different addressees and actors must not be merged into a single plural command to dwell.
% mdparagraph:P285:end

\end{minipage}\par\addvspace{4pt}

\par\begin{minipage}{\linewidth}\raggedright
% mdparagraph:P286:begin
\subsection*{Cook 48.9}
% mdparagraph:P286:end

% key: KEY-EX-XII-XIII-48-09
% mdparagraph:P287:begin
Dōð swā, ðæt gē ne forweorðen.
% mdparagraph:P287:end

% output: EX-XII-XIII-48-09-O001
% mdparagraph:P288:begin
\textbf{Cook’s prompt:} Do thus, that ye perish not.
% mdparagraph:P288:end

% mdparagraph:P289:begin
\textbf{Target and explanation:} \textbf{Dōð swā} is a plural command, “do thus.” The following \textbf{ðæt gē ne forweorðen} gives its negative purpose, so that they may not perish, with second-plural present subjunctive. The target is an adverb of manner: \textbf{swā} or \textbf{ðus}, not the object pronoun \textbf{ðis} “this.” An indicative statement that they are already perishing reverses the purpose.
% mdparagraph:P289:end

\end{minipage}\par\addvspace{4pt}

\par\begin{minipage}{\linewidth}\raggedright
% mdparagraph:P290:begin
\subsection*{Cook 48.10}
% mdparagraph:P290:end

% key: KEY-EX-XII-XIII-48-10
% mdparagraph:P291:begin
Nū gē gesēoð hū hit mid mē is, and gē hīerað hwæt ic tō ēow sprece.
% mdparagraph:P291:end

% output: EX-XII-XIII-48-10-O001
% mdparagraph:P292:begin
\textbf{Cook’s prompt:} Now ye see how it is with me, and ye hear what I say (sprecan) to (tō) you.
% mdparagraph:P292:end

% mdparagraph:P293:begin
\textbf{Target and explanation:} Both main verbs have plural \textbf{gē}: you see the stated circumstances and hear the speech. \textbf{Hū hit mid mē is} is the supplied state-of-affairs expression “how it is with me,” not mere physical accompaniment. \textbf{Hwæt ic tō ēow sprece} is what Joseph says to them: first-singular speaker, plural addressees. Equivalent English wording such as “how matters stand with me” is acceptable; replacing either embedded clause by a new direct question loses the sentence’s structure.
% mdparagraph:P293:end

\end{minipage}\par\addvspace{4pt}

\reserve{7\baselineskip}
% mdparagraph:P294:begin
\section*{R. The princess’s gifts and the banquet’s end}
% mdparagraph:P294:end

% mdparagraph:P295:begin
The assigned reading is Cook 175.18–177.2, ending with \textbf{geręstan} before \emph{Apollonius as Teacher}. The divisions below cover every clause. Page and line references retain the edition’s pages and the stated physical-line convention. Use the source notes and complete local help throughout.
% mdparagraph:P295:end

\reserve{8\baselineskip}
% mdparagraph:P296:begin
\subsection*{R01}
% mdparagraph:P296:end

% key: KEY-NEW-U18-R01
% mdparagraph:P297:begin
\textbf{Task:} Translate the complete assigned passage in the following three consecutive divisions. All local help and original notes remain available.
% mdparagraph:P297:end

% output: NEW-U18-R01-O001
\reserve{5\baselineskip}
% mdparagraph:P298:begin
\textbf{Translate 175.18–176.2, through goldhorde.}
% mdparagraph:P298:end

% mdparagraph:P299:begin
Indeed, when the king’s daughter saw that Apollonius was so well trained in all good accomplishments, her heart fell in love with him. Then, after the end of the banquet, the maiden said to the king: ‘Dear father, a little while ago you allowed me to give Apollonius whatever I wished from your treasure.’
% mdparagraph:P299:end

% mdparagraph:P300:begin
\textbf{Target and explanation:} The daughter is the perceiver and the person who falls in love; Apollonius is the person admired. “Accomplishments,” “skills” and “arts” can render the supplied cræft sense here. The permission belongs to an earlier occasion, recalled with “a little while ago”; it does not command a gift at that earlier moment. The free choice and royal treasure remain part of what she recalls.
% mdparagraph:P300:end

% output: NEW-U18-R01-O002
\reserve{5\baselineskip}
% mdparagraph:P301:begin
\textbf{Translate 176.3–176.13, through gesāwon.}
% mdparagraph:P301:end

% mdparagraph:P302:begin
King Arcestrates said to her: ‘Give him whatever you wish.’
% mdparagraph:P302:end

% mdparagraph:P303:begin
She then went out very gladly and said: ‘Master Apollonius, with my father’s permission I give you two hundred pounds of gold, and four hundred pounds by weight of silver, and the greatest portion of precious clothing, and twenty male slaves.’ And she then spoke thus to the slaves: ‘Carry these things with you, which I promised Apollonius my teacher, and put them in the chamber in front of my friends.’
% mdparagraph:P303:end

% mdparagraph:P304:begin
This was then done in this way, according to the princess’s command; and all the men who saw them praised her gifts.
% mdparagraph:P304:end

% mdparagraph:P305:begin
\textbf{Target and explanation:} Retain the two separate amounts and their materials: two hundred pounds of gold, four hundred pounds by weight of silver. \textbf{Ðone mǣstan dǣl} is a superlative: “the greatest part/\allowbreak{}portion” or “most of the precious clothing,” not merely “a large amount.” The twenty male slaves are also a gift, then become the addressees of the carrying and placing commands. “Teacher” or “master” can render Lārēow; chamber is supplied with the original dining-room note. The princess’s commands are then reported as carried out. The final relative means that the men saw the gifts; original note 13 expressly supplies męnn as its antecedent.
% mdparagraph:P305:end

% output: NEW-U18-R01-O003
\reserve{5\baselineskip}
% mdparagraph:P306:begin
\textbf{Translate from Ðā sōðlice geęndode in 176.13 through 177.2 geręstan.}
% mdparagraph:P306:end

% mdparagraph:P307:begin
Then the banquet ended, and all the men rose and took leave of the king and the princess, wished them well, and went home. Apollonius likewise said: ‘You good king and one who has mercy on the poor, and you princess, lover of learning, may you be well.’
% mdparagraph:P307:end

% mdparagraph:P308:begin
He also looked toward the slaves whom the maiden had given him, and said to them: ‘Take with you these things which the princess gave me, and let us go and seek our lodging, so that we may rest.’
% mdparagraph:P308:end

% mdparagraph:P309:begin
\textbf{Target and explanation:} The guests’ rising, leave-taking and going home are narrated as completed. “Bade them farewell” can render the wish that they be well; grētton is appropriate to taking leave here. Apollonius first addresses the king and princess, then turns to the slaves given him. His last words are commands and a proposal, followed by the purpose of rest. They do not narrate arrival or sleep. “Inn,” “guesthouse” and “lodging” can render giesthūs in this context. The printed cwēn is the princess, the established king’s daughter, rather than a newly introduced queen.
% mdparagraph:P309:end

\reserve{8\baselineskip}
% mdparagraph:P310:begin
\subsection*{R02}
% mdparagraph:P310:end

% key: KEY-NEW-U18-R02
% mdparagraph:P311:begin
\textbf{Task:} For each cited form, name the referent and explain its grammatical role. Keep the whole passage available.
% mdparagraph:P311:end

% output: NEW-U18-R02-O001
\reserve{5\baselineskip}
% mdparagraph:P312:begin
\textbf{Explain hiere and his in 175.20.}
% mdparagraph:P312:end

% mdparagraph:P313:begin
\textbf{Hiere} is genitive feminine singular, the daughter as possessor of \textbf{mōd}. \textbf{His} is genitive masculine singular referring to Apollonius; in the supplied love expression the affection is directed toward him. Thus “her heart fell in love with him,” or idiomatically “she fell in love with him.” It does not say that he fell in love with her.
% mdparagraph:P313:end

% output: NEW-U18-R02-O002
\reserve{5\baselineskip}
% mdparagraph:P314:begin
\textbf{Identify the speaker and addressee of ic and ðē in 176.5.}
% mdparagraph:P314:end

% mdparagraph:P315:begin
The princess is \textbf{ic}, the nominative giver and speaker of this new quotation. Apollonius is \textbf{ðē}, dative recipient of \textbf{giefe}, “I give you.” The preceding king’s speech has ended; keeping him as speaker would assign the gifts to the wrong person.
% mdparagraph:P315:end

% output: NEW-U18-R02-O003
\reserve{5\baselineskip}
% mdparagraph:P316:begin
\textbf{Distinguish him in 176.19 from him in 176.20.}
% mdparagraph:P316:end

% mdparagraph:P317:begin
In \textbf{ðe ðæt mǣden him forgiefen hæfde}, him is dative singular Apollonius, recipient of the slaves. In \textbf{and him cwæð tō}, him is dative plural, the slaves whom he now addresses; tō is postposed. The identical spellings have different numbers and referents, established by their clauses and the following plural commands.
% mdparagraph:P317:end

\reserve{8\baselineskip}
% mdparagraph:P318:begin
\subsection*{R03}
% mdparagraph:P318:end

% key: KEY-NEW-U18-R03
% mdparagraph:P319:begin
\textbf{Task:} Parse the following selected source expressions and explain their functions. Supplied lexical and quantity explanations remain visible.
% mdparagraph:P319:end

% output: NEW-U18-R03-O001
\reserve{5\baselineskip}
% mdparagraph:P320:begin
\textbf{Compare the grammatical subjects of geseah 175.18 and gefēoll 175.20.}
% mdparagraph:P320:end

% mdparagraph:P321:begin
The subject of \textbf{geseah} is \textbf{ðæs cyninges dohtor}, the king’s daughter. The grammatical subject of \textbf{gefēoll} is singular \textbf{hiere mōd}, her heart or mind; hiere alone is a genitive possessor. Idiomatic English may say “she fell in love,” but that does not change the Old English subject’s structure.
% mdparagraph:P321:end

% output: NEW-U18-R03-O002
\reserve{5\baselineskip}
% mdparagraph:P322:begin
\textbf{Explain the relationships expressed by punda and goldes in 176.6.}
% mdparagraph:P322:end

% mdparagraph:P323:begin
\textbf{Punda} is genitive plural of the counted weight-unit after \textbf{twā hund}. \textbf{Goldes} is genitive singular identifying the material measured. The supplied whole phrase means “two hundred pounds of gold.” Calling both simply “possessors” obscures the different dependencies; this task applies the visible quantity analysis rather than testing an unassigned numeral system.
% mdparagraph:P323:end

% output: NEW-U18-R03-O003
\reserve{5\baselineskip}
% mdparagraph:P324:begin
\textbf{Use original note 13 to explain ðe hīe gesāwon in 176.13 and the role of giefa.}
% mdparagraph:P324:end

% mdparagraph:P325:begin
Original note 13 expressly identifies \textbf{męnn} as the antecedent of \textbf{ðe}. The relative represents those men as the subject of \textbf{gesāwon}; \textbf{hīe} is accusative plural, the gifts they saw. \textbf{Giefa} is accusative plural object of \textbf{hęredon}, with genitive hiere naming the princess as their giver/\allowbreak{}possessor. Thus “all the men who saw them praised her gifts.” The intervening words do not make gifts the relative subject.
% mdparagraph:P325:end

% output: NEW-U18-R03-O004
\reserve{5\baselineskip}
% mdparagraph:P326:begin
\textbf{Compare the case, number and agreement of gesunde in 176.16 and 176.18.}
% mdparagraph:P326:end

% mdparagraph:P327:begin
In \textbf{bǣdon hīe gesunde bēon} (176.16), \textbf{gesunde} is accusative plural, agreeing with accusative \textbf{hīe}, the king and princess as subject of infinitive bēon. In \textbf{bēon gē gesunde} (176.18), it is nominative plural, agreeing with \textbf{gē} in a direct second-plural subjunctive wish. Both formulas mean “be well” or “farewell.” The adjective is not an adverb merely because English says “well”; the two bēon forms also have different grammatical functions.
% mdparagraph:P327:end

\reserve{8\baselineskip}
% mdparagraph:P328:begin
\subsection*{R04}
% mdparagraph:P328:end

% key: KEY-NEW-U18-R04
% mdparagraph:P329:begin
\textbf{Task:} Describe the clause relations in the following spans using the complete source and its local support.
% mdparagraph:P329:end

% output: NEW-U18-R04-O001
\reserve{5\baselineskip}
% mdparagraph:P330:begin
\textbf{Analyse the opening sentence 175.18–20: identify the temporal frame, the seeing-content clause and the main event.}
% mdparagraph:P330:end

% mdparagraph:P331:begin
\textbf{Mid-ðȳ-ðe} introduces the temporal frame, “when the king’s daughter saw ….” Within it, \textbf{ðæt Apollonius … wæs getogen} is the content seen: his education and accomplishments. \textbf{Ðā gefēoll hiere mōd …} gives the ensuing main event. Geseah and wæs are indicative in this asserted circumstance. These are nested temporal and content relations followed by a main clause, not three interchangeable relative clauses.
% mdparagraph:P331:end

% output: NEW-U18-R04-O002
\reserve{5\baselineskip}
% mdparagraph:P332:begin
\textbf{Compare ðæt ic mōste 176.1 with swā-hwæt-swā 176.2 and 176.3–4.}
% mdparagraph:P332:end

% mdparagraph:P333:begin
\textbf{Ðæt ic mōste giefan …} is the supplied permission-content construction: “you allowed me to give …,” as original note 2 directs. Mōste expresses permission here, not obligation. \textbf{Swā-hwæt-swā} means “whatever,” the freely selected thing given: the daughter recalls what she \textbf{wolde}, and the king now says she may give what she \textbf{wille}. These are free-choice relatives, not direct questions. The earlier wolde is first-singular past subjunctive; wille is second-singular present subjunctive addressed to her.
% mdparagraph:P333:end

% output: NEW-U18-R04-O003
\reserve{5\baselineskip}
% mdparagraph:P334:begin
\textbf{Analyse the final ðæt clause 177.2 and explain how it relates to gān wē sēcean.}
% mdparagraph:P334:end

% mdparagraph:P335:begin
\textbf{Gān wē} is a first-plural present-subjunctive exhortation, “let us go.” Plain infinitive \textbf{sēcean} gives the action sought by going, finding their lodging. \textbf{Ðæt wē mægen ūs geręstan} gives its purpose, “so that we may rest”: mægen is first-plural present subjunctive, geręstan an infinitive, and \textbf{ūs} the accusative reflexive directed by original note 177.2 (§184b). Neither the exhortation nor the purpose narrates an arrival or completed rest.
% mdparagraph:P335:end

\reserve{8\baselineskip}
% mdparagraph:P336:begin
\subsection*{R05}
% mdparagraph:P336:end

% key: KEY-NEW-U18-R05
% mdparagraph:P337:begin
\textbf{Task:} Explain the passage’s sequence from the actual statements, distinguishing completed actions from proposals.
% mdparagraph:P337:end

% output: NEW-U18-R05-O001
\reserve{5\baselineskip}
% mdparagraph:P338:begin
\textbf{What authority does the princess give for her gifts, and who selects them?}
% mdparagraph:P338:end

% mdparagraph:P339:begin
She invokes her father’s earlier permission, then announces the gifts \textbf{be mīnes fæder lēafe}, with his permission. He confirms that she may give whatever she wishes. The daughter selects and enumerates the gifts; the king does not individually prescribe the stated amounts, clothing and slaves.
% mdparagraph:P339:end

% output: NEW-U18-R05-O002
\reserve{5\baselineskip}
% mdparagraph:P340:begin
\textbf{What does the final speech establish about Apollonius’s next action, and what remains unnarrated at this boundary?}
% mdparagraph:P340:end

% mdparagraph:P341:begin
He orders the slaves to take the gifts with them and proposes that they go to seek their lodging so that they can rest. The earlier guests have already gone home, but this last speech does not narrate Apollonius’s arrival, sleep or completed rest. “He plans/\allowbreak{}proposes to go and rest” is accurate; “he went to sleep” imports an event beyond the selected boundary.
% mdparagraph:P341:end

\reserve{7\baselineskip}
% mdparagraph:P342:begin
\subsection*{Cook’s original reading notes}
% mdparagraph:P342:end

% mdparagraph:P343:begin
These twenty-two notes retain Cook’s wording and all twenty-nine anchor relationships. Their section numbers refer to his grammar; the visible modern help supplies the interpretations used above. His etymological questions and the Beowulf comparison are retained apparatus, not additional exercises or a poetry assignment.
% mdparagraph:P343:end

% mdparagraph:P344:begin
\textbf{175, note 10} (175.19 getogen,). See getēon. What relation has getogen to Mod. Eng. wanton?
% mdparagraph:P344:end

% mdparagraph:P345:begin
\textbf{175, note 11} (175.21 “Lēofa). See 55.
% mdparagraph:P345:end

% mdparagraph:P346:begin
\textbf{176, note 1} (176.01 lȳtle; 176.01 ǣr,). Lat. paulo ante. See 178.
% mdparagraph:P346:end

% mdparagraph:P347:begin
\textbf{176, note 2} (176.01 ðæt; 176.01 ic; 176.01 mōste). Translate by the infinitive sign, to. The OE. follows the Latin.
% mdparagraph:P347:end

% mdparagraph:P348:begin
\textbf{176, note 3} (176.04 wille.). See 197.
% mdparagraph:P348:end

% mdparagraph:P349:begin
\textbf{176, note 4} (176.04 ūtēode,). Not in Latin.
% mdparagraph:P349:end

% mdparagraph:P350:begin
\textbf{176, note 5} (176.05 “Lārēow). Lat. magister.
% mdparagraph:P350:end

% mdparagraph:P351:begin
\textbf{176, note 6} (176.06 punda; 176.06 punda). See 154. c.
% mdparagraph:P351:end

% mdparagraph:P352:begin
\textbf{176, note 7} (176.06 goldes,; 176.07 seolfres,). See 153. f.
% mdparagraph:P352:end

% mdparagraph:P353:begin
\textbf{176, note 8} (176.07 gewihte). See 174.
% mdparagraph:P353:end

% mdparagraph:P354:begin
\textbf{176, note 9} (176.07 dēorwurðes). MS. deorwurðan.
% mdparagraph:P354:end

% mdparagraph:P355:begin
\textbf{176, note 10} (176.08 ðēowa; 176.08 manna.). Lat. servos.
% mdparagraph:P355:end

% mdparagraph:P356:begin
\textbf{176, note 11} (176.11 būre). Lat. triclinio.
% mdparagraph:P356:end

% mdparagraph:P357:begin
\textbf{176, note 12} (176.12 hǣse). See the derivation of Mod. Eng. behest.
% mdparagraph:P357:end

% mdparagraph:P358:begin
\textbf{176, note 13} (176.13 ðe). Refers to męnn.
% mdparagraph:P358:end

% mdparagraph:P359:begin
\textbf{176, note 14} (176.14 ārison,). So in Beowulf (653–655): “Werod eall ārās; grētte ðā . . . guma ōðerne, . . . and him hǣl ābēad.”
% mdparagraph:P359:end

% mdparagraph:P360:begin
\textbf{176, note 15} (176.16 bēon,). Lat. vale dicentes.
% mdparagraph:P360:end

% mdparagraph:P361:begin
\textbf{176, note 16} (176.17 earmra; 176.18 lāre). See 153. d.
% mdparagraph:P361:end

% mdparagraph:P362:begin
\textbf{176, note 17} (176.18 gesunde.). Lat. valete.
% mdparagraph:P362:end

% mdparagraph:P363:begin
\textbf{176, note 18} (176.20 hæfde,). See 188.
% mdparagraph:P363:end

% mdparagraph:P364:begin
\textbf{177, note 1} (177.01 gān). See 193. a.
% mdparagraph:P364:end

% mdparagraph:P365:begin
\textbf{177, note 2} (177.02 ūs). See 184. b.
% mdparagraph:P365:end

% mdparagraph:P366:begin
\textsuperscript{1}: Cook, \emph{Exercises in Old English} (1895), pp. 29–33, joint XII, XIII, Exercises 44–48; \emph{A First Book in Old English}, second edition (1895), pp. 175–177 and the cited grammatical sections. The command-passive reference in 45.5 uses 177.10 as a grammatical parallel without assigning the following episode. Additional checked lexical evidence is Bosworth–Toller, main volume pp. 95 and 1056 and Supplement pp. 81, 240 and 646; proper-name and supplied nominal evidence is documented in the unit’s source apparatus. The earlier findan form funde is explicitly allowed by Cook’s strong-verb discussion; no correction to its earlier teaching is required.
% mdparagraph:P366:end

\end{document}
